\documentclass[11pt]{article}

\usepackage[T1]{fontenc}
\usepackage[utf8]{inputenc}
\usepackage[english]{babel}
\usepackage[margin=1in]{geometry}
\usepackage{lmodern}
\usepackage{microtype}
\usepackage{mathtools}
\usepackage{amsmath,amssymb,amsthm}
\usepackage{array}
\usepackage{booktabs}
\usepackage{tabularx}
\usepackage{float}
\usepackage{graphicx}
\usepackage{tikz}
\usepackage{caption}
\usepackage{enumitem}
\usepackage{fancyhdr}
\usepackage[numbers,sort&compress]{natbib}
\usepackage{doi}
\usepackage{hyperref}
\usepackage{cleveref}
\usepackage{orcidlink}
\usepackage{titling}
\makeatletter
\g@addto@macro\UrlBreaks{\do\_\do\.\do\/\do\-}
\makeatother

 
\newtheorem{theorem}{Theorem}[section]
\newtheorem{lemma}[theorem]{Lemma}
\newtheorem{proposition}[theorem]{Proposition}
\newtheorem{corollary}[theorem]{Corollary}
\theoremstyle{definition}
\newtheorem{definition}[theorem]{Definition}
\newtheorem{remark}[theorem]{Remark}
\floatstyle{ruled}
\newfloat{algorithm}{tbp}{loa}
\floatname{algorithm}{Algorithm}

\urlstyle{tt}
\captionsetup{
  font=small,
  labelfont=bf,
  labelsep=period
}
\captionsetup[table]{skip=6pt,position=top}
\captionsetup[figure]{skip=6pt}
\usetikzlibrary{arrows.meta,fit,positioning,shapes.geometric}
\crefname{algorithm}{algorithm}{algorithms}
\Crefname{algorithm}{Algorithm}{Algorithms}
\crefname{theorem}{theorem}{theorems}
\Crefname{theorem}{Theorem}{Theorems}
\crefname{lemma}{lemma}{lemmas}
\Crefname{lemma}{Lemma}{Lemmas}
\crefname{proposition}{proposition}{propositions}
\Crefname{proposition}{Proposition}{Propositions}
\crefname{corollary}{corollary}{corollaries}
\Crefname{corollary}{Corollary}{Corollaries}
\crefname{definition}{definition}{definitions}
\Crefname{definition}{Definition}{Definitions}
\crefname{remark}{remark}{remarks}
\Crefname{remark}{Remark}{Remarks}
\setlist[enumerate]{leftmargin=*, itemsep=2pt, topsep=4pt}
\setlength{\droptitle}{-3.2em}
\pretitle{\begin{center}\LARGE\bfseries}
\posttitle{\par\end{center}\vspace{0.5em}}
\preauthor{\begin{center}\normalsize}
\postauthor{\par\end{center}\vspace{-0.4em}}
\predate{\begin{center}\small}
\postdate{\par\end{center}\vspace{0.6em}}
\setlength{\emergencystretch}{2em}
\clubpenalty=10000
\widowpenalty=10000
\displaywidowpenalty=10000
\interfootnotelinepenalty=10000
\allowdisplaybreaks[2]
\fancypagestyle{firstpage}{\fancyhf{}
  \fancyfoot[C]{\scriptsize Automorph Inc., Wilmington, DE, USA \quad \texttt{ioannis@automorph.io} \quad \textcopyright\ 2026 \quad CC-BY 4.0}
  \renewcommand{\headrulewidth}{0pt}
  \renewcommand{\footrulewidth}{0pt}
}

\title{Carrier Exactification after Promotion on Fixed Support in Six Birds Theory}
\author{Ioannis Tsiokos\,\orcidlink{0009-0009-7659-5964}}
\date{\today}
\hypersetup{
  pdftitle={Carrier Exactification after Promotion on Fixed Support in Six Birds Theory},
  pdfauthor={Ioannis Tsiokos},
  hidelinks
}

\begin{document}

\maketitle
\thispagestyle{firstpage}

\begin{abstract}
Earlier papers in the fixed-support line of Six Birds Theory established admissible event packages, predictive and recombination quotients, a control-stable Route~A flagship corpus, promotion criteria, and the intrinsic promoted-package survivor. They left open the internal carrier status of that survivor. We introduce a carrier-level substrate on the same Route~A flagship corpus, consisting of three state-like objects and one structural relation. The three objects are a \emph{lawful same-branch witness}, a refined witness rung on which the same-branch transport and readout law exactifies; a \emph{coarse public signature shadow}, which retains public signature role but not lawful transport; and a \emph{flagship-only protocol-indexed carrier surface}, an autonomous lift of the witness on which protocol coordinates are first-class state data. The structural relation is \emph{non-faithful public compression}: a projection from the witness to the public shadow that preserves public signature role but does not preserve the refined lawful structure. A seven-criterion exactification audit adapted from standard operational-theory methodology---readout sufficiency, transport and readout law, non-reducibility against the lower quotient baselines, section and protocol honesty, denominator and expansion stress, and cross-seed support---together with a refreshed section/protocol honesty rerun on the lifted flagship patch and a controlled adjacent-surface search stabilises the atlas. Three claim-bearing quadratic families and two proto-event/test families survive on the scoped carrier. Carrier law on the public surface, cross-seed carrier exactification, and full-core carrier exactification are not earned; those non-claims are preserved explicitly. A Lean formal layer records the criterion vocabulary and a typed mirror of the final verdict. The carrier atlas, with its scope and compression structure made explicit, is the substrate carried forward.
\end{abstract}

\medskip
\noindent\textbf{Keywords:} carrier exactification; promoted package; fixed support; lawful witness; protocol-indexed carrier; Six Birds Theory.

\section{Introduction}
\label{sec:introduction}

The fixed-support line of Six Birds Theory has progressed in stages. The opening paper established admissible event packages and the discipline of comparing candidate descriptions on a common support. The second introduced current and predictive quotients on single histories and showed that future-relevant structure can survive below the current-event interface. The third lifted the same discipline to branch assemblages and showed that branchwise predictive content and later recombination-sensitive content can separate on fixed support. The fourth isolated a small control-stable Route~A flagship family with a theorem-backed core. The fifth asked whether that witness family supports a genuine promoted package and showed that at least one intrinsic candidate survives on fixed support. The sixth then analyzed the internal structure of that survivor and left open a still sharper question: what kind of promoted object, if any, has actually survived \citep{Tsiokos2026EventPackages,Tsiokos2026HolonomyMemory,Tsiokos2026MarkingErasureRecombination,Tsiokos2026RecombinationWitnesses,Tsiokos2026PromotionCriteria,Tsiokos2026InternalStructure}.

The present paper asks whether the promoted survivor exactifies into a \emph{carrier} rather than remaining only a descriptor partition, a process family, a chart-only presentation, or a witness surface. The substrate we contribute is an answer to that question: on the same Route~A flagship corpus, a carrier-level atlas consisting of three state-like objects and one structural relation between them.

\paragraph{Contributed substrate.} The three objects are:
\begin{itemize}[leftmargin=2em, itemsep=1pt, topsep=2pt]
\item the \emph{lawful same-branch witness}, a refined witness rung on which the same-branch transport and readout law exactifies;
\item the \emph{coarse public signature shadow}, which retains public signature and readout role but does not carry the refined lawful transport; and
\item the \emph{flagship-only protocol-indexed carrier surface}, an autonomous lift of the witness on which protocol coordinates are first-class state data.
\end{itemize}
The structural relation is \emph{non-faithful public compression}: a projection from the witness to the public shadow that preserves public signature role without preserving the refined lawful structure. Together, these four items form the carrier-level atlas the series carries forward.

\paragraph{How the atlas is earned.} A seven-criterion exactification audit adapted from standard operational-theory methodology---readout sufficiency, transport and readout law, non-reducibility against the lower quotient baselines inherited from earlier papers, section and protocol honesty, denominator and expansion stress, and cross-seed support---isolates which candidate objects survive as carrier-level rungs. A refreshed section-and-protocol honesty rerun on the lifted flagship patch and a controlled adjacent-surface search stabilise the scope restriction. The witness passes all seven criteria on the flagship core. The public shadow passes only signature-level surface. The protocol-indexed carrier passes at flagship-only indexed scope. The audit is adapted, not new; the substrate it stabilises is.

Two downstream diagnostics---three claim-bearing quadratic families and two proto-event/test families---survive on the flagship-only protocol-indexed carrier. Carrier law on the public surface, a cross-seed carrier, and a full-core carrier are not earned; cross-seed canonicalization remains partial and four nearby positives remain real scope limitations. A narrow Lean formal layer records the criterion vocabulary and a typed mirror of the final verdict, without claiming to derive the empirical verdict inside Lean.

\paragraph{Scope.} We do not claim carrier law at the public surface, cross-seed carrier exactification, full-core carrier exactification, nonindexed promotion of the scoped lifted surface, global quadratic carrying, or a general event geometry. The result is narrower: the promoted package exactifies into a lawful same-branch witness together with a flagship-only protocol-indexed carrier surface; the public carrier remains a signature-level shadow, and the stronger carrier lanes do not land.

The remainder of the paper is organized as follows. \Cref{sec:paper6-starting-point} recalls the Paper~6 starting surface. \Cref{sec:candidate-classes-reductions-criteria} fixes the candidate classes, reduction targets, and exactification criteria. \Cref{sec:corpus-artifacts-pipeline} records the corpus, artifact surface, and pipeline discipline. \Cref{sec:live-surface-exactification} presents the live-surface exactification attempt. \Cref{sec:same-branch-law-and-public-shadow} develops the lawful same-branch witness and the public shadow split. \Cref{sec:protocol-indexed-flagship-carrier} presents the scoped protocol-indexed carrier surface. \Cref{sec:cross-seed-reach-and-scope-contracts} and \Cref{sec:stress-expansion-honesty-new-surface} collect the cross-seed, scope, stress, honesty, and new-surface limitations. \Cref{sec:quadratic-carrying-and-proto-event-pilot} records the scoped quadratic and proto-event diagnostics. \Cref{sec:final-gate-and-carrier-atlas} gives the final gate and carrier atlas, \Cref{sec:formal-and-lean-criterion-layer} summarizes the formal layer, \Cref{sec:discussion} interprets the result and records the limitations, \Cref{sec:philosophical-reading} gives a short philosophical reading of the witness/public-shadow split, and \Cref{sec:conclusion} states the final claim.
 \section{Starting point: promoted packages after Paper 6}
\label{sec:paper6-starting-point}

Papers~1--6 progressively narrowed the admissible comparison base, clarified
what kinds of exact-finite structure survive below current readout, and isolated a
small Route~A corpus on which promoted-package questions can be asked under matched
controls. This section recalls what the earlier papers earned, what the Paper~6
survivor is, and why that survivor still leaves a carrier question open
\citep{Tsiokos2026EventPackages,Tsiokos2026HolonomyMemory,Tsiokos2026MarkingErasureRecombination,Tsiokos2026RecombinationWitnesses,Tsiokos2026PromotionCriteria,Tsiokos2026InternalStructure}.

\subsection{Fixed-support packages and global obstruction}

Paper~1 formulated the fixed-support event-package problem
\citep{Tsiokos2026EventPackages}. Local contexts can be internally coherent, Boolean,
and mutually compatible in the small, yet still fail to assemble into a single
admissible global package on the same support. The basic question is not whether one
can invent a convenient summary for each local context in isolation, but whether a
summary survives once all comparisons are forced onto a common finite substrate.

Three consequences follow. First, fixed support forbids silent movement of the
comparison base: a putative structural gain must survive on the same corpus rather than
being purchased by changing interfaces or enlarging the underlying sample of histories.
Second, local admissibility does not automatically determine a lawful or globally
legible object. Third, the evidence surface is exact-finite; the core claims are tied
to concrete surfaces on which admissibility, obstruction, and separation can be audited
exactly.

The classical marginal problem \citep{Fine1982}, the Kochen--Specker theorem
\citep{KochenSpecker1967}, and the sheaf-theoretic formulation of non-locality and
contextuality \citep{AbramskyBrandenburger2011} are the nearest prior analogues
of this local-to-global discipline.

A promoted package may therefore survive as a meaningful object in one sense and still
fall short of carrier status in a stronger sense. Paper~7 asks whether the promoted
survivor from the later papers crosses that stronger boundary.

\subsection{Predictive quotients and recombination witnesses}

Papers~2 and~3 show that fixed support can carry structure not exhausted by
current-visible summaries
\citep{Tsiokos2026HolonomyMemory,Tsiokos2026MarkingErasureRecombination}. Paper~2
introduced the contrast between a current quotient and a predictive quotient on single
histories: two histories can agree at the current interface while still differing in
future-relevant structure. The most publicly legible surface need not be the surface on
which the decisive law is carried. Computational-mechanics causal states
\citep{ShaliziCrutchfield2001}, information-bottleneck predictive summaries
\citep{TishbyPereiraBialek1999}, and the behavioral-systems emphasis on observed
trajectories \citep{Willems1991} are the nearest trajectory-level analogues.

Paper~3 lifted that contrast from single histories to branch assemblages, separating
branchwise predictive structure from recombination-sensitive structure and showing that
the latter can refine the former on fixed support. The consistent-histories
\citep{Griffiths1984} and decoherent-histories \citep{GellMannHartle1990} programmes
supply the nearest prior language for assemblage-level quotient structure. A current
summary, a predictive quotient, and a richer recombination-sensitive quotient can
coexist without collapsing into the same object.

If future-relevant and recombination-sensitive structure can persist below current
readout, then a public summary need not be the right place to look for the strongest
lawful object. Papers~2--3 did not themselves produce a carrier, but they supplied the
exact-finite logic for asking whether an object with more law than the public readout
can be earned on the same frozen support.

\subsection{The Route A flagship corpus}

Paper~4 narrowed the empirical surface to the Route~A flagship family and its matched
comparisons \citep{Tsiokos2026RecombinationWitnesses}. The series fixed a small
control-stable corpus with two flagship anchors, a local neighborhood of nearby
positives, required controls, and a formal kernel aligned to the flagship surface.
Paper~7 inherits this corpus.

Two consequences follow. The flagship anchors provide a stable decision surface against
which survival, separation, and honesty are read. And nearby positives and controls
keep their distinct roles: nearby positives remain informative local comparisons rather
than theorem-bearing anchors by proximity, and controls remain structural discriminators
rather than decorative nulls.

Paper~7 therefore does not ask whether some broader search surface contains a
carrier-like description somewhere. It asks what, if anything, exactifies on the frozen
Route~A surface with the promoted package from Papers~5--6 as the starting point.

\subsection{Promotion and the Paper 6 survivor}

Paper~5 posed the next sharper question \citep{Tsiokos2026PromotionCriteria}: does the
Route~A witness family support a genuine promoted package, or only lower-layer witness
structure? Its answer was positive but carefully delimited.
On the frozen Route~A corpus, at least one promoted package survives. The promoted
package result was not a claim about arbitrary higher-level objecthood; it was a
criterion-bound survival statement on a fixed comparison surface with explicit controls
and admissibility conditions.

Paper~6 moved from existence to internal structure
\citep{Tsiokos2026InternalStructure}. In its audit, the intrinsic route emerged as the
main survivor: the package carried descended readout support, signature-level update or
closure support, second-stage stability on the decisive surface, retained invariant and
quadratic families, local chart or readiness signals, and protocol-honesty structure at
the package level. The survivor was not a vacuous label pasted on top of lower-layer
evidence.

The status earned in Paper~6 nevertheless remained package-level. Descended readout
support is not yet a public carrier law; signature-level update support is not yet an
autonomous transport law; retained quadratic families and local chart signals do not
yet amount to a global event geometry. Paper~6 left something more precise than a
surviving name but less than a carrier verdict.

The series had not yet determined whether the surviving object should be understood as
only a promoted package, as a lawful same-branch witness, as a publicly legible
carrier, as a protocol-indexed surface, or as something reducible to a lower baseline.
Package survival and internal structure were necessary preparation for the carrier
question, but not its answer.

\subsection{The remaining carrier gap}

A surviving promoted package is not automatically a carrier. Carrier language requires
more than package survival, signature-level closure cues, or local readiness signals:
one needs to know on what surface transport and readout law are actually supported,
whether the relevant object is public or only refined, whether it survives honesty and
reduction audits, and whether its scope extends beyond the flagship-local surface on
which it was first isolated.

Paper~6 did not identify carrier law on the coarse readout surface, did not establish
an autonomous carrier law on a designated state surface independent of the later
exactification audit, did not show that the surviving structure extends across the
broader frozen comparison surface as a uniform carrier, and did not justify treating
retained quadratic or local chart signals as if they already defined a full
higher-level geometry.

Paper~7 therefore asks: given the promoted package that survives after Paper~6, does
any candidate exactify into a carrier, and if so on what surface and with what scope?
The question forces three distinctions that the rest of the paper keeps separate: the
coarse public surface, the lawful witness surface, and any stronger carrier surface
that survives only under additional indexing or scope restriction. The next section
makes those distinctions operational by fixing the candidate classes, reduction
baselines, and exactification criteria.
 \section{Carrier candidates, reductions, and exactification criteria}
\label{sec:candidate-classes-reductions-criteria}

Paper~6 established that a promoted package survives and carries meaningful internal
structure, but did not settle what kind of object the survivor is. Some candidates are
only descriptors of lower behavior; some are process families; some are chart-like
local representations; some are genuine carrier candidates; and some survive only under
explicit indexing or scope restrictions. This section fixes that vocabulary, states the
reduction baselines against which richer objects are tested, and records the
exactification criteria used in the later audit.

\subsection{Object types after promotion}

A promoted object may still be only a descriptor of lower behavior, only a process
family, only a local chart, or only a protocol-conditioned representation. These
object types must be separated before any candidate is evaluated against the
exactification criteria. The carrier language we use is deliberately weaker than
classical objecthood: we do not assume that a promoted package corresponds to a
traditional "thing" in the sense rejected by ontic structural realism
\citep{LadymanRoss2007}, nor that emergent higher-level carriers are reducible to
their microphysical substrate \citep{Anderson1972}.

At the weakest end are \emph{descriptor partitions}. A descriptor partition groups lower
states or histories into equivalence classes that summarize visible or predictive
behavior. It is a legitimate object for comparison, but it does not by itself provide a
new transport or readout law. In the present paper, descriptor partitions include the
current-visible and predictive quotients inherited from Paper~2 and the branchwise and
recombination quotients inherited from Paper~3
\citep{Tsiokos2026HolonomyMemory,Tsiokos2026MarkingErasureRecombination}. Their role in
ticket31 is baseline-level: later candidates must do more than merely rename one of
these already-available summaries.

Next come \emph{process object families}. A process family can organize explicit runtime
states, histories, or action surfaces without thereby becoming a promoted carrier. The
series already uses such process-like layers to generate benchmark behavior and to build
the quotients against which later claims are tested. Process families may have their own
transitions and observables, but those transitions remain bookkeeping or executable
runtime structure unless they descend to a stable promoted state space with an honest
readout and transport role of their own.

A third object type is the \emph{chart-only intrinsic object}: a local patch or
coordinate-style representation that fits a restricted region of the flagship surface
without surviving changes of section, protocol, or surface. We use chart language in
the standard differential-geometric sense \citep{Lee2013}: a chart fits a neighborhood
of a manifold without necessarily gluing to a global atlas. Chart-only objects can
look geometrically or algebraically rich while remaining strictly local, which is why
C3 cannot be reduced to a vague novelty condition.

Beyond chart-only objects lie \emph{chart plus section/fiber families}. These candidates
attach explicit section choices or fiber data to the promoted base state, in the
sense of sheaf-theoretic sections over a base space \citep{MacLaneMoerdijk1992}. They
organize readout and update over a broader representation family, but they risk
depending on representational choice rather than earning a section-honest state law. Section
dependence may be a legitimate part of the object, or it may signal that the apparent
object is only a charted presentation of something simpler.

\emph{Protocol-indexed objects} carry explicit protocol or internalization data as part
of the state description. A route that survives only after protocol state has been
internalized is not automatically spurious, but it is not the same as a nonindexed
carrier. An indexed object may survive honestly as indexed while still falling short of
a context-free or public carrier claim.

\emph{Carrier candidates} are state-like objects intended to support both readout and
transport or update on their own state space. The Lean layer expresses this abstractly
as a designated surface together with a candidate state space, a readout relation, and
a transport relation. A carrier candidate may be intrinsic, section-mediated,
fiber-mediated, or indexed, and is evaluated by whether it carries an honest
state/readout/transport role rather than only a descriptive or charting role.

An \emph{exactified carrier} is a carrier candidate that survives the full
exactification audit: sufficient for the intended readout panel, supporting transport
and readout law, noncollapsing into simpler baselines, honest under section and protocol
changes, stress- and expansion-robust, and extending beyond the flagship pair when a
broader claim is at stake. Candidates that fail some gates may remain positive in a
weaker sense without being exactified.

Table~\ref{tab:candidate-taxonomy} summarizes the candidate classes, the reduction or
fragility pressure each faces, and the range of statuses each can receive. Later
sections determine which candidate classes survive which parts of the audit.

\begin{table}[t]
\centering
\small
\caption{Candidate taxonomy used in the carrier-exactification audit. The table
separates descriptor, process, chart, indexed, and carrier lanes before any final gate
is assigned.}
\label{tab:candidate-taxonomy}
\begin{tabularx}{\linewidth}{@{}>{\raggedright\arraybackslash}p{0.20\linewidth}>{\raggedright\arraybackslash}p{0.20\linewidth}>{\raggedright\arraybackslash}p{0.22\linewidth}>{\raggedright\arraybackslash}X@{}}
\toprule
Candidate class & Role in the audit & Main risk & Possible final status \\
\midrule
Descriptor partition & baseline quotient layer & richer notation merely re-encodes $Q$, $M$, $K$, $R$, or $\Lambda^\ast$ & descriptor/process only \\
Process object family & executable runtime competitor & runtime bookkeeping mistaken for promoted state law & descriptor/process only \\
Chart-only intrinsic & local patch representation & local fit fails section, protocol, or surface transport & chart-only survivor \\
Chart plus section/fiber & section-mediated object family & section dependence mistaken for objecthood & unresolved; chart-like; or live candidate \\
Protocol-indexed object & indexed state/readout family & indexed survival restyled as context-free object & scoped or indexed-only survival \\
Intrinsic carrier candidate & base promoted-state carrier lane & collapse to simpler witness or baseline surface & live candidate; scoped carrier; or exactified carrier \\
Section/fiber carrier candidate & bundle/fiber carrier lane & fiber structure collapses to chart or descriptor role & live candidate; scoped carrier; or unresolved \\
New-surface carrier candidate & prospective adjacent-surface lane & no materialized surface or executable law yet & prospective or no exactification \\
\bottomrule
\end{tabularx}
\end{table}

\subsection{Reduction baselines}

The candidate taxonomy becomes operational only once richer objects are compared against
explicit simpler targets. The fixed-support series already provides those targets. Paper~2
introduced the current-visible quotient $Q$ and the predictive quotient $M$, together
with the factor map
\[
\pi : M \to Q,
\]
which forgets future-relevant structure when that structure is invisible at the current
surface \citep{Tsiokos2026HolonomyMemory}. Paper~3 introduced the branchwise quotient
$K$ and the recombination quotient $R$, together with the forgetful map
\[
\eta : R \to K,
\]
which collapses recombination-sensitive distinctions back to branchwise predictive
structure when no stricter recombination refinement survives
\citep{Tsiokos2026MarkingErasureRecombination}. Ticket31 then adds one more baseline
target, $\Lambda^\ast$, the minimal behavioral carrier or reachable-image baseline: the
smallest already-available law-like image that can inherit the observable and transition
behavior without introducing a richer promoted object.

These baselines are the reduction targets used in the candidate census and reduction
map. A candidate is \emph{reducible} when its states, readouts, and transport or update
roles factor through one of these simpler targets, or through an already-landed simpler
positive lane, without losing the claim one is trying to make. A reducible candidate
may still be a useful presentation, but it has not contributed a new state space
together with a new readout/transport law.

This is the operational content of C3. Non-reducibility does not mean ``visibly more
complicated'' or ``uses a richer notation''; it means that the candidate does something
that cannot be supplied by the simpler baseline or already-earned lane for the claim
under consideration. A section-mediated or protocol-indexed object may look richer than
$Q$, $M$, $K$, $R$, or $\Lambda^\ast$ and still fail C3 if its apparent law factors
through those targets. A stronger-looking carrier surface may fail the strongest-carrier
claim if it contributes no state, readout, or transport role beyond an already-landed
lawful witness or witness/public-shadow basis. C3 blocks the relabeling of simpler
structure as a stronger object; it does not reward novelty for its own sake. In this
respect C3 is only loosely analogous to a weak-emergence screen
\citep{Bedau1997}: for the claim at hand, the candidate must contribute behavior not
already recoverable from the simpler baseline.

Descriptor, process, and chart lanes therefore remain in the paper as explicit
comparison targets rather than as failed drafts of the same object. Section~\ref{sec:final-gate-and-carrier-atlas}
returns to this point when the final gate separates strongest exactification claims
from weaker survived-positive surfaces.

\subsection{Exactification criteria}

Carrier exactification is governed by seven criteria. Together they answer a single
question: does the candidate support an honest state/readout/transport role strong
enough for the claim being made? The criteria are cumulative but not uniformly binding:
failure of a broader criterion such as cross-seed support can block a strongest claim
without erasing a narrower positive surface. Operational reconstructions of physical
theories in terms of such criterion bundles \citep{Hardy2001,ChiribellaDArianoPerinotti2011}
are the nearest prior methodological analogue.

\begin{enumerate}[label=\textbf{C\arabic*.},leftmargin=*,itemsep=0.6em]
\item \textbf{Readout sufficiency.} The candidate must support the intended readout panel
on the surface for which the claim is made. In exact-finite terms, the designated seeds
must admit valid candidate states with honest readout. Failure of C1 means the proposed
object is not even readable on the target surface.

\item \textbf{Transport or readout law.} The candidate must support a coherent transport,
update, or closure role together with readout compatibility on its own state space.
Failure of C2 means that the object remains descriptive or signature-like rather than
lawful as a carrier candidate.

\item \textbf{Non-reducibility.} The candidate must not collapse to $Q$, $M$, $K$, $R$,
$\Lambda^\ast$, or another already-landed simpler lane for the relevant claim. Failure
of C3 does not say the candidate is useless; it says the candidate does not justify a
stronger object-status label than the simpler surface already provides.

\item \textbf{Section honesty.} The candidate must survive admissible section or
completion changes on the surface relevant to the claim. Failure of C4 means the
apparent object is tied to a section choice or chart patch that the audit is not
prepared to treat as honest object data.

\item \textbf{Protocol honesty.} The candidate must survive explicit protocol
internalization. A route that works only because omitted protocol state has been hidden
in the presentation fails C5. By contrast, an object that survives only after protocol
state has been internalized may still earn a \emph{protocol-indexed} status. What it
does not earn automatically is a nonindexed carrier verdict.

\item \textbf{Stress and expansion honesty.} The candidate must remain intact under the
denominator-stress and repo-native surface-expansion panels, without treating branch
breaks, collapses, or expansion artifacts as wins. Failure of C6 blocks stronger
carrier claims even when a narrower positive lane remains.

\item \textbf{Cross-seed support.} Any cross-seed or full-core claim must survive beyond
the two flagship anchors. In the present fixed corpus, that means canonicalization to
nearby positives or a comparably strong extension surface. Failure of C7 blocks
cross-seed or full-core exactification, but it does not by itself erase a flagship-local
or indexed positive survivor.
\end{enumerate}

Two asymmetries should be noted. First, failure of public law does not erase a lawful
witness: the coarse public surface and a lawful witness surface answer different
questions and can separate. Second, failure of C7 does not erase all positivity: a
candidate may still survive honestly as flagship-only or protocol-indexed even when
the broader extension needed for cross-seed or full-core language does not materialize.
Without these asymmetries, every shortfall of the strongest claim would collapse into
a purely negative verdict.

\subsection{Scoped survival versus full exactification}

The exactification audit is not all-or-nothing. The final gate uses the criteria above
to distinguish several kinds of positive and negative outcome.

At the strongest end lies \emph{full carrier exactification}: a candidate supports
readout and transport law, survives the honesty and stress panels, remains
nonreducible, and extends far enough across the frozen comparison surface to warrant
cross-seed or full-core carrier language.

Below it lies \emph{flagship-only carrier exactification}: a carrier candidate is
honest and nontrivial on the flagship surface but does not earn a broader extension. A
distinct possibility is \emph{protocol-indexed carrier-only survival}, where the
candidate remains a real positive surface but only with explicit protocol indexing on
a limited scope.

Lower on the ladder are the \emph{lawful witness} and the \emph{public signature
shadow}. A lawful witness is already a real positive object but need not coincide with
the coarse public surface. A public signature shadow may support readout-facing
legibility without carrying the same law as the witness. Witness law, public shadow,
and scoped carrier survival are kept distinct rather than merged.

Finally, \emph{new-surface exactification} resolves the carrier question only after
moving to an adjacent surface.

The formal meaning of these categories is fixed here; the empirical allocation of
candidates to them is deferred to the later audit and final-gate sections.

The Lean layer mirrors this vocabulary at the criterion level. It formalizes designated
surfaces, carrier candidates, reductions, honesty predicates, exactification evidence,
and weaker scoped or protocol-indexed evidence, together with small abstract gating
lemmas. That formal layer stabilizes the logic of the criteria, but it does not certify
the empirical ticket31 artifacts or replay the final gate from repository execution.
 \section{Corpus, baselines, and pipeline}
\label{sec:corpus-artifacts-pipeline}

The audit is performed on a frozen corpus rather than on an open search surface. We ask
what exactifies on the Route~A surface inherited from Papers~4--6, not whether one can
enlarge the search until some stronger object appears. The audit fixes one finite
corpus, resolves the comparison baselines on that corpus, and records the evidence
through machine-readable artifacts and a serial canonical pipeline. Detailed source
maps, command surfaces, and asset-index entries are in
\Cref{app:reproducibility-and-source-map}.

\subsection{The frozen seed corpus}

The frozen corpus contains eleven seeds. Their role counts are fixed in advance and do
not change later in the paper: two flagship anchors, four nearby positives, four
controls, and one support-only row. The flagship anchors are
\texttt{t27.flagship.reference\_positive} and
\texttt{t27.flagship.completion\_partner}. They form the primary carrier-exactification
anchor panel inherited from the Route~A flagship family of Paper~4
\citep{Tsiokos2026RecombinationWitnesses}. The nearby positives are the four
relative-cycle shortlist cases \texttt{t27.nearby.n3\_d1\_w1of3\_2of3} and
\texttt{t27.nearby.n4\_d1\_w1of3\_2of3}, together with
\texttt{t27.nearby.n5\_d1\_w1of2\_1of2} and
\texttt{t27.nearby.n5\_d2\_w1of3\_2of3}. They are not coequal anchors. Their role is to
test whether a candidate that looks strong on the two-anchor core extends beyond that
core in a way that could matter for cross-seed or broader-scope claims.

The four controls are the null and memory controls
\texttt{t27.control.null\_baseline} and \texttt{t27.control.memory\_only}. The
remaining two are the artifact and boundary controls
\texttt{t27.control.artifact} and \texttt{t27.control.boundary\_collapse}.
Together they test the main alternative explanations already stabilized in the earlier
series: trivial mixture, predictive memory without the stronger target structure,
artifact or omitted-state dependence, and boundary-sensitive collapse. The remaining row,
\texttt{t27.support.erasure\_recovery\_anchor}, is retained as support-only. It is part
of the family context and remains scientifically useful, but it is not an exact-readout
or lawful-transport anchor and is not promoted to ordinary carrier evidence.

The corpus is a fixed comparison surface with a small positive core, a nearby challenge
panel, a matched control bundle, and one support-only row. It is the surface on which
the candidate classes of Section~\ref{sec:candidate-classes-reductions-criteria} are
tested; later sections distinguish flagship, nearby, control, and support-only
outcomes on this same frozen corpus rather than by changing the comparison base.
Table~\ref{tab:ticket31-corpus-overview} summarizes the corpus and its handling rules.

\begin{table}[t]
\centering
\small
\caption{Frozen ticket31 corpus and panel handling. The table separates ordinary direct
seeds, the paired artifact control, and the support-only row before any later verdict is
assigned.}
\label{tab:ticket31-corpus-overview}
\begin{tabularx}{\linewidth}{@{}>{\raggedright\arraybackslash}p{0.18\linewidth}>{\raggedright\arraybackslash}p{0.12\linewidth}>{\raggedright\arraybackslash}p{0.24\linewidth}>{\raggedright\arraybackslash}p{0.20\linewidth}>{\raggedright\arraybackslash}X@{}}
\toprule
Role group & Count & Function in the audit & Handling mode & Main later interpretation \\
\midrule
Flagship anchors & 2 & primary carrier/law anchor panel & ordinary direct core & main positive support surface \\
Nearby positives & 4 & cross-seed and non-flagship challenge panel & ordinary, non-anchor seeds & later scope test rather than automatic transfer \\
Null, memory, boundary controls & 3 & null and alternative-explanation checks & ordinary controls & contrast surfaces, not positive transfer cases \\
Artifact control & 1 logical seed & artifact and internalization check & paired logical panel & not ordinary positive support; remains separate throughout \\
Support-only erasure row & 1 & family-context and erasure-side support & support-only panel & not exact readout or lawful-transport transfer \\
\bottomrule
\end{tabularx}
\end{table}

\subsection{Panels and scope conventions}

The paper uses several overlapping but nonidentical panels. The \emph{direct core panel}
is the ordinary seed panel on which candidates are first compared without paired-control
special handling: the two flagship anchors together with the ordinary controls that can
be treated as single resolved seeds. This is the panel on which the most basic
candidate-versus-control questions are first asked. The \emph{paired artifact panel} is
kept separate from that direct core. The artifact control is represented by two protocol
endpoints, not by a single ordinary row, because the whole point of that control is to
test whether apparent structure survives explicit internalization rather than hiding it.

The \emph{nearby-positive panel} consists of the four nearby positives. These are not
used as silent extensions of the flagship core. They are the non-flagship panel against
which broader scope claims are challenged. The \emph{support-only panel} contains the
erasure-recovery anchor and is treated as support-only throughout: useful for family
context and interpretation, but not counted as exact readout or law transfer. Together,
the nearby-positive and support-only panels form the main non-flagship comparison
surface. Later sections use that surface to distinguish real scope limitation,
support-only transfer, and ordinary flagship evidence rather than collapsing them into a
single broad notion of positivity.

Controls are not positive transfer cases; the support-only row is not an exact-readout
row in disguise; and the paired artifact endpoints are not ordinary seeds that can be
averaged back into the main count. Fixing these conventions before the result sections
prevents later positive-surface language from silently widening the claim base.

\subsection{Baseline resolution and reduction adapters}

The corpus is frozen, but the baseline layer is not inherited as a complete frozen
payload: many of the expected ticket30 derived artifacts were not materialized in the
live repository snapshot. The baseline is therefore built by \emph{source recovery}.
The seed corpus, surface registry, and baseline smoke artifacts recover the relevant
benchmark cases, control semantics, and comparison surfaces from repo-native code,
tests, and committed benchmark configurations, rather than assuming that old frozen
JSON outputs are already present.

In practice, baseline resolution uses three modes. Six ordinary seeds are resolved by
\emph{direct benchmark configuration}: the two flagship anchors, three ordinary controls,
and the support-only row all point to committed benchmark configs that can be reloaded
directly. The four nearby positives are resolved as \emph{synthetic relative-cycle
cases}: they are not invented ad hoc, but reconstructed from the relative-cycle family
parameters that define the nearby shortlist. The artifact control is resolved as a
\emph{logical paired control}: two protocol endpoints are kept together as one control
object rather than flattened into a single ordinary seed. The reduction-baseline smoke
surface records this split explicitly. Ten rows are fully usable as ordinary executable
baseline bundles, while the artifact row remains paired-only by design.

These resolved seeds are then passed through the reduction adapters introduced in
Section~\ref{sec:candidate-classes-reductions-criteria} and inherited from the quotient
layers of Papers~2 and~3 \citep{Tsiokos2026HolonomyMemory,Tsiokos2026MarkingErasureRecombination}.
For each executable seed, the
baseline layer computes the current quotient $Q$, predictive quotient $M$, factor map
$\pi : M \to Q$, branchwise quotient $K$, recombination quotient $R$, factor map
$\eta : R \to K$, and minimal behavioral carrier baseline $\Lambda^\ast$. For the paired
artifact control, the same baseline family is built memberwise and then retained as a
paired comparison object. This is the concrete way in which the C3 reduction targets are
instantiated in ticket31: they are not abstract ideals floating above the corpus, but
computed or reconstructed baseline bundles attached to the resolved seeds.

The same source-recovery discipline governs the operative surfaces. Exact-readout,
support-only, signature-only operator, section-family, protocol, primitive-action, and
prospective new-surface families are recorded in the surface registry together with
their validation basis. Some of these surfaces are builder-only or source-defined
rather than inherited as old frozen outputs; the main claim is tied to the repo-native
surface as rebuilt and indexed here, not to an imagined complete archive.

\subsection{Artifact-backed evidence model}

Each major step in the audit writes machine-readable summaries or registries. The
frozen seed corpus, the operative surface registry, the candidate census, and the
reduction map establish what is compared and which lanes are treated as baseline,
descriptor/process, chart-only, live carrier candidates, or unresolved. Later steps
write more specialized summaries for live-surface exactification, law and witness
audits, scope contracts, stress and honesty panels, downstream scoped diagnostics, and
the final gate.

Operational notes record why a builder or gate was run, how paired or support-only
handling was interpreted, and what caveats must remain visible. The main text does not
reproduce those notes as logs but depends on their decisions; figures and tables
synthesize the artifact-backed results. The claims come from a controlled artifact
surface rather than from a retrospective narrative summary.

\subsection{Canonical rebuild and reproducibility}

The paper's reproducibility story is intentionally narrow and explicit. The audit does
not rely on an informal sequence of shell commands run by hand. Instead, a canonical
runner and an asset index record the executable or validated path through the audit.
The runner is one-command and explicitly serial. Serial execution matters because
several builders and phase runners rewrite shared derived JSON; parallel execution would
change the meaning of the evidence surface rather than merely speeding it up. The
canonical phase order therefore moves from seed freeze and baseline recovery, through
live-surface exactification, PICA-backed law and bridge recovery, the strengthening lane,
the cross-seed, scope-contract, falsifier, and final-gate decision lane, Lean
witness-bridge validation, and finally asset indexing and contract checks.

The runner also makes the external dependencies visible. PICA and vendor-backed cargo
steps are either executed or explicitly validated as part of the canonical surface, not
left to manual reconstruction outside the pipeline. The Lean criterion layer and witness
bridge are likewise included in the validation path. At the same time, the reproducible
surface remains honest about its own limits. The asset index records known gaps, marks
environment-dependent artifacts such as run manifests or some cargo-backed bundle logs,
and distinguishes rebuildable artifacts from validated-only surfaces. In the current
ticket31 status, the canonical surface is ready with known gaps and validated in
\texttt{check-only} mode; the paper uses that indexed surface rather than pretending that
every historical comparison lane was freshly rebuilt in the present turn.

Reproducibility includes claim-regression discipline. The canonical runner checks that
the final contract remains scoped: the final gate may not silently widen into
cross-seed, full-core, public-side law, or broader downstream claims; the public
shadow must remain distinct from the lawful witness; the paired artifact panel must
remain separate; and the support-only row must remain support-only. Our canonical
runner, asset index, and artifact-backed evidence model follow the computational
reproducibility discipline advocated by Donoho \citep{Donoho2010} and the FAIR
data principles \citep{WilkinsonEtAl2016}. Detailed commands, file-level source maps,
and asset-hash information are in \Cref{app:reproducibility-and-source-map}. The
evidence is frozen, serially reconstructible, and audited against regressions in both
artifacts and claim language.

\begin{figure}[t]
\centering
\resizebox{0.94\linewidth}{!}{\begin{tikzpicture}[
  font=\small,
  node distance=1.05cm and 2.6cm,
  stage/.style={
    draw,
    rounded corners,
    align=center,
    fill=black!4,
    text width=0.36\linewidth,
    minimum height=1.15cm,
    inner sep=5pt
  },
  side/.style={
    draw,
    rounded corners,
    align=center,
    fill=black!2,
    text width=0.24\linewidth,
    minimum height=1.1cm,
    inner sep=5pt
  },
  flow/.style={-Latex, thick},
  relation/.style={-Latex, thick, dashed}
]
\node[stage] (paper6) {Paper 6 promoted package};
\node[stage, below=of paper6] (census) {Candidate census and reduction audit};
\node[stage, below=of census] (live) {Live-surface carrier candidate};
\node[stage, below=of live] (witness) {Lawful same-branch witness};
\node[side, left=of witness] (public) {Public signature shadow\\\texttt{signature\_only}};
\node[side, right=of witness] (indexed) {Scoped protocol-indexed carrier\\\texttt{flagship\_only}};
\node[stage, below=of indexed] (gates) {Cross-seed, stress, honesty,\\and new-surface gates};
\node[stage, below=of gates] (verdict) {Final verdict\\\texttt{protocol\_indexed\_carrier\_only}};

\draw[flow] (paper6) -- (census);
\draw[flow] (census) -- (live);
\draw[flow] (live) -- (witness);
\draw[relation] (witness) -- node[above, align=center, font=\scriptsize, fill=white, inner sep=1pt] {non-faithful\\public compression} (public);
\draw[flow] (witness) -- (indexed);
\draw[flow] (indexed) -- (gates);
\draw[flow] (gates) -- (verdict);
\end{tikzpicture}
 }
\caption{Ticket31 audit path from the Paper~6 promoted package to the final gate. The
path isolates the lawful same-branch witness, keeps the public row as a signature-level
shadow, and records the final verdict only after the later scope, stress, honesty, and
adjacent-surface gates have been applied.}
\label{fig:ticket31-pipeline}
\end{figure}

Figure~\ref{fig:ticket31-pipeline} gives the main flow of the audit. The diagram ends
at the scoped verdict actually earned and does not imply cross-seed or full-core
exactification.
 \section{Live-Surface Exactification}
\label{sec:live-surface-exactification}

This section asks whether the surviving promoted package can be materialized as a
live-surface carrier candidate before any later law extraction, same-branch witness
refinement, or protocol-lift strengthening. The question is narrower than the final
gate: it is not whether the object carries public-surface law, a cross-seed carrier,
or a lawful same-branch witness, but whether the intrinsic promoted object can be
organized as a coherent state surface with nontrivial readout and transport roles on
the direct core panel, with the paired artifact control kept explicit and without
collapse into the descriptor or process baselines of Section~\ref{sec:candidate-classes-reductions-criteria}.

The answer is positive but preliminary. An object-construction step produces a concrete
intrinsic live-surface object proposal; an exactification audit on the direct core panel,
together with the paired artifact panel, shows that the intrinsic branch is a real carrier
candidate rather than merely a descriptor/process reencoding or a chart-only presentation;
a comparison-branch audit then weighs the same intrinsic object against the main alternative
live-surface lanes. That comparison keeps the bundle/section-fiber route visible as a
secondary carrier-style lane, classifies the sectioned route as chart-only, and records
that an early protocol-indexed interface does not yet instantiate coherently on the direct
core panel. The present section reports that early live-surface verdict only. The stronger
same-branch lawful witness and the later scoped protocol-indexed carrier surface are
developed in \Cref{sec:same-branch-law-and-public-shadow,sec:protocol-indexed-flagship-carrier}.

\subsection{Constructing the intrinsic live-surface object}

The object-construction step builds a serializable object proposal on the frozen
corpus from the intrinsic promoted survivor. The resulting candidate, the
\texttt{intrinsic\_naked\_interface}, is a \emph{proposal}: the step shows that the
promoted package can be organized as a state-like surface suitable for exactification
tests, not that the later law and honesty gates have been passed.

The proposal is built in a seed-tagged patch representation. Its states are identified
as \texttt{seed\_tagged\_patch\_state}: precise enough to serialize the flagship
anchors, the ordinary controls, and the paired artifact endpoints into one direct
comparison object, but not yet advertised as a seed-independent or globally canonical
state law. The construction records seven states. Five correspond to the ordinary
direct-core seeds: the two flagship anchors together with the null, memory, and
boundary controls. The remaining two are the uninternalized and internalized artifact
endpoints, kept separate as a logical paired panel in accordance with the paired
control rule from Section~\ref{sec:corpus-artifacts-pipeline}; the artifact row is not
flattened into one ordinary seed.

The proposal also records eight readout entries and three transport-family entries.
There is a concrete readout family on the flagship anchors and control rows, and a
transport family attached to the same object. The transport semantics are, however,
\emph{signature-only}: the proposal does not yet claim explicit target closure, lawful
composition, or autonomous transport law. It provides the minimal live-surface
structure needed for the next gate: a finite state proposal whose readout and update
surfaces can be tested on the direct core panel while the paired artifact control is
preserved as a separate logical panel.

The construction does not prove that the intrinsic branch is a final carrier. It shows
that the Paper~6 intrinsic survivor can be recast as a coherent live-surface object
proposal rather than an informal internal-structure bundle. The proposal is
patch-scoped, seed-tagged, and transport-limited, but already more than a descriptor or
chart sketch, and can be pressed against readout, transport, and reduction pressure as
a carrier candidate.

\subsection{Readout, transport, and reduction pressure on the core panel}

The first exactification audit then evaluates that intrinsic live-surface proposal. The
evaluation scope is the ordinary direct core together with the explicit paired artifact
panel. On the ordinary direct core, the candidate is tested on five executable seeds:
the two flagship anchors and the three ordinary controls. The artifact control remains
outside that direct-core count and continues to be tracked as a logical paired panel.
This preserves the panel discipline established in Section~\ref{sec:corpus-artifacts-pipeline}
while allowing the intrinsic branch to be judged on a genuine direct-core comparison
surface.

The intrinsic branch supports readout on all five ordinary direct-core seeds, with two
exact cases and three partial cases. Exact and partial support appear on the same
coherent state proposal rather than on unrelated local summaries, the direct core is
not compressed to a singleton, the two flagships are separated from the controls, and
state/readout assignments are tied to the intrinsic patch surface. The proposal is
therefore readable on the intended early panel and is not merely a descriptor
partition.

The transport result is weaker. Transport remains signature-level only: no explicit
closed transport law is claimed, and no composition test is recorded. The section does
not yet speak of a lawful same-branch witness; that witness must be earned on the
stronger audit of \Cref{sec:same-branch-law-and-public-shadow}. The intrinsic live
surface supports nontrivial transport-side structure without collapsing into a null
bookkeeping layer, but signature support is present and explicit law is not.

The reduction picture is mixed. On the raw seed-tagged view, the intrinsic branch
remains noncollapsed under the current proxy checks. On the retagged view, pressure
rises and seed-tagging dependence is high, but there is still no explicit baseline
collapse. The reduction pressure does not justify demoting the branch to a pure $Q$-,
$M$-, $K$-, $R$-, or $\Lambda^\ast$-style baseline presentation, so the audit lands at
\texttt{carrier\_candidate} rather than chart-only, descriptor/process only, or
collapse.

Two clarifications prevent overclaim. The result is strictly live-surface and
direct-core: it says nothing yet about nearby positives, cross-seed canonicalization,
or full-core extension. And signature closure is not explicit law, so the result does
not license lawful transport. The intrinsic branch emerges as a serious carrier
candidate---nontrivial, readable on the direct core, transport-bearing in a signature
sense, and not explicitly collapsed---but not as the final lawful object.

\subsection{Comparison branches}

The comparison-branch audit then asks whether the intrinsic branch is genuinely the best
live-surface route or only one representative among several equally plausible alternatives.
The comparison is run on the same direct-core plus paired-artifact panel, so the branch
verdicts are directly comparable rather than being artifacts of a changed panel.

The intrinsic branch remains the best live-surface branch. The comparison artifact
records the reason as matching the strongest live-surface status while keeping the
least mediated object surface and the richest readout panel. The intrinsic candidate
does not require an explicit section choice or fiber presentation to instantiate on the
direct core, which makes it the cleanest early carrier lane before any stronger law
has been extracted.

The main non-intrinsic competitor is the \texttt{bundle\_section\_fiber\_interface},
which also lands as a \texttt{carrier\_candidate}. It matches the intrinsic branch on
the ordinary direct core: all five direct-core seeds are supported, two exactly and
three partially, and no explicit retagged collapse is recorded. It loses the
comparison not because it is empty but because it is more mediated; its fiber-tagged
state space makes it a less direct live-surface object. The branch remains a secondary
carrier-style lane rather than the best early branch.

The \texttt{sectioned\_package\_interface} shows positive support on the direct core
with the same two exact and three partial seed outcomes but is classified as
\texttt{chart\_only}. Readout support and noncollapsed proxy pressure are not by
themselves enough to earn carrier language; the sectioned branch remains tied to
section-mediated presentation rather than an object-level state law, and so remains a
chart-plus-section family.

The early \texttt{protocol\_indexed\_interface} is weaker still, supporting only the
two flagship seeds on the direct core with no exact cases and no coherent control-side
coverage; its status is \texttt{failed\_to\_instantiate}. Protocol-index internalization
at this stage remains reference-only and does not cover the full direct core
coherently. The later scoped protocol-indexed survivor of
\Cref{sec:protocol-indexed-flagship-carrier} is a different construction with its own
honesty and promotion gates, not a relabeling of this early interface.

Table~\ref{tab:live-surface-summary} summarizes the live-surface results before any later
law extraction or protocol lift is introduced.

\begin{table}[t]
\centering
\footnotesize
\caption{Live-surface exactification and comparison results before the later
law-extraction stage. The intrinsic branch is the best live-surface carrier candidate,
but transport is still signature-level at this stage.}
\label{tab:live-surface-summary}
\begin{tabularx}{\linewidth}{@{}>{\raggedright\arraybackslash}p{0.16\linewidth}>{\raggedright\arraybackslash}p{0.13\linewidth}>{\raggedright\arraybackslash}p{0.17\linewidth}>{\raggedright\arraybackslash}p{0.15\linewidth}>{\raggedright\arraybackslash}p{0.16\linewidth}>{\raggedright\arraybackslash}X@{}}
\toprule
Stage / branch & Scope & Readout result & Transport result & Reduction / mediation issue & Status \\
\midrule
Intrinsic object construction & direct core plus paired artifact panel & 8 source-recovered readouts on a 7-state seed-tagged patch proposal & 3 transport-family entries; signature-only semantics & patch-scoped, seed-tagged construction; paired artifact kept logical & proposal \\
Intrinsic exactification audit & 5 ordinary direct-core seeds plus separate paired artifact panel & 5/5 supported, with 2 exact and 3 partial outcomes & signature-level support only; no explicit law or composition claim & raw pressure noncollapsed; retagged pressure higher; no explicit baseline collapse & carrier candidate \\
Sectioned comparison branch & same panel & direct-core support present but section-mediated & signature-only and section-mediated & positive support remains chart-like rather than object-level & chart only \\
Bundle/section-fiber comparison branch & same panel & matches the intrinsic core-support pattern on the direct core & signature-only carrier-style support & more mediated by section/fiber structure than intrinsic & carrier candidate \\
Protocol-indexed comparison branch & same panel & only partial flagship coverage; no coherent control-side instantiation & signature-only reference rows do not rescue the lane & reference-only protocol internalization on the early panel & failed instantiation \\
\bottomrule
\end{tabularx}
\end{table}

\subsection{Live-surface verdict}

The live-surface stage ends with a bounded positive result. The audit does not remain
at the level of descriptor partitions, process families, or chart-only presentations;
it produces an intrinsic live-surface carrier candidate that is state-bearing,
readable on the direct core, paired-panel aware, and not explicitly collapsed by the
reduction audit.

The result ceiling is nevertheless low. The intrinsic branch is still patch-scoped and
seed-tagged, transport remains signature-level rather than lawful, and no cross-seed
or public-surface-law claim is in play. The correct verdict is therefore not ``carrier
exactified'' but ``best live-surface carrier candidate.''

The next section asks whether that best live-surface branch supports a genuine
same-branch law, and how its lawful witness separates from the coarse public shadow.
 \section{Same-Branch Law and the Public Signature Shadow}
\label{sec:same-branch-law-and-public-shadow}

Section~\ref{sec:live-surface-exactification} showed that the intrinsic branch is a
live-surface carrier candidate with signature-level transport on the early direct-core
panel. The PICA-backed law audit now asks where law actually lands on that same
branch. The audit uses a fixed same-branch witness ladder rather than a single
undifferentiated object, and the outcome is the central structural split of the paper:
law exactifies on a refined same-branch witness, while the coarse public carrier
remains a public signature shadow. We identify the rung on which the same-branch law
becomes lawful, and explain why that law does not transfer to the coarse public
carrier under enrichment, authentication, or upgrade attempts.

\subsection{The same-branch law ladder}

The PICA-backed law audit runs on a fixed five-candidate same-branch ladder whose
entries move from the public-facing coarse layer to progressively refined witness
surfaces:
\begin{enumerate}[label=\arabic*.,leftmargin=2.0em,itemsep=0.2em]
\item \texttt{intrinsic\_naked\_interface\_coarse},
\item \texttt{intrinsic\_naked\_interface\_witness\_panel\_refined},
\item \texttt{intrinsic\_naked\_interface\_witness\_observable\_enriched},
\item \texttt{intrinsic\_naked\_interface\_witness\_deeper\_object\_map},
\item \texttt{intrinsic\_naked\_interface\_witness\_rewrite\_effective\_operator}.
\end{enumerate}
The ladder prevents loose talk of ``the'' carrier law where the repository tests
several nested same-branch candidates with different statuses.

The status pattern is decisive. The coarse public candidate and the panel-refined
witness candidate both remain \texttt{signature\_only}. The observable-enriched and
deeper-object-map candidates collapse at the law stage: they do not sustain a coherent
law after retagging collapse. Only the fifth rung,
\nolinkurl{intrinsic_naked_interface_witness_rewrite_effective_operator}, lands as
\texttt{lawful}, and is selected as the first lawful candidate in the ladder. This is
the lawful same-branch witness.

The corrected non-PICA rerun does not give a lawful object; it leaves the primary
candidate at \texttt{signature\_only}. The PICA-backed law layer upgrades the
same-branch law verdict, but only at the refined witness rung, not at the coarse
public rung. The upgrade is real but selective.

The ladder yields three conclusions. The early live-surface candidate of
\Cref{sec:live-surface-exactification} is not illusory: a refined same-branch witness
emerges as lawful. Refining the public panel or enriching observables is not by itself
sufficient; two intermediate witness-like rungs still collapse. And the coarse public
carrier does not inherit law simply because law appears somewhere on the same branch;
the successful law lives on a refined witness surface with its own
state/readout/transport role.

The lawful witness closes transport with explicit or mixed agreement and retains
stable readout after retagging collapse, whereas the coarse public candidate remains
coherent only at the signature level. That is the law-theoretic difference between
the two surfaces.

\subsection{The lawful witness and the public shadow}

The main structural split can now be stated plainly. The lawful same-branch witness is
\nolinkurl{intrinsic_naked_interface_witness_rewrite_effective_operator}; the coarse
public carrier is \nolinkurl{intrinsic_naked_interface_coarse} with status
\texttt{signature\_only}; and the relation between them is recorded as
\nolinkurl{non_faithful_public_compression}.

The coarse public carrier is retained as a public signature or readout shadow: the
public-facing layer is coherent enough to preserve a stable signature surface, which
matters for the later operational basis of the audit. The public shadow does not,
however, carry the same transport law as the refined witness, and merging the two
levels would erase the law-audit result.

The phrase \emph{non-faithful public compression} names the relation recorded by the
artifacts: a projection from the lawful witness surface to the public surface exists,
but it does not preserve the full lawful structure. The information-theoretic analogue
is a lossy compression in the sense of rate--distortion theory \citep{CoverThomas2006},
and the structural analogue is the ontic-vs-epistemic distinction for quantum states
\citep{HarriganSpekkens2010}: public-surface features need not carry all of the
structure present at the refined level. What survives on the public side is public
signature or readout role, not the refined same-branch law. In the final atlas, the
public entry is therefore classified as a public-facing shadow while the same-branch
witness is the lawful survivor.

Calling the public object a ``failed carrier'' would collapse two distinct statements
into one negative slogan. The same-branch law exists but exactifies on a refined
witness; the coarse public carrier remains coherent only at the signature level. The
split is evidence about where law lives, not evidence that the public surface is
meaningless.

The final gate and carrier atlas continue to track the lawful same-branch witness and
the public shadow as separate surfaces, and the cutover basis keeps the
witness/public-shadow pair explicit rather than promoting the coarse public layer into
a lawful carrier.

\begin{figure}[t]
\centering
\begin{tikzpicture}[
  font=\small,
  node distance=2.2cm,
  layer/.style={
    draw,
    rounded corners,
    align=center,
    fill=black!4,
    text width=0.34\linewidth,
    minimum height=1.55cm,
    inner sep=6pt
  },
  relation/.style={-Latex, thick, dashed}
]
\node[layer] (witness) {Lawful same-branch witness\\[2pt]
  \scriptsize \texttt{intrinsic\_naked\_interface\_witness\_}\\
  \scriptsize \texttt{rewrite\_effective\_operator}\\[2pt]
  lawful};
\node[layer, right=of witness, fill=black!2] (public) {Public signature shadow\\[2pt]
  \scriptsize \texttt{intrinsic\_naked\_interface\_coarse}\\[2pt]
  \texttt{signature\_only}};

\draw[relation] (witness) -- node[above, align=center, font=\scriptsize, fill=white, inner sep=1pt] {non-faithful\\public compression} (public);
\end{tikzpicture}
 \caption{Two-level witness/public split in schematic form. The lawful same-branch witness
and the coarse public signature shadow remain distinct surfaces, and their relation is
non-faithful public compression rather than lawful transport on the public side.}
\label{fig:witness-public-shadow}
\end{figure}

Figure~\ref{fig:witness-public-shadow} gives the schematic split, while
Table~\ref{tab:witness-public-shadow-split} records the same relation with the explicit
ticket31 row names and statuses.

\begin{table}[t]
\centering
\footnotesize
\caption{Two-level law/readout split. The law exactifies on the refined same-branch
witness, while the public carrier remains a signature-level shadow related by
non-faithful public compression.}
\label{tab:witness-public-shadow-split}
\begin{tabularx}{\linewidth}{@{}>{\raggedright\arraybackslash}p{0.20\linewidth}>{\raggedright\arraybackslash}p{0.30\linewidth}>{\raggedright\arraybackslash}p{0.14\linewidth}>{\raggedright\arraybackslash}p{0.18\linewidth}>{\raggedright\arraybackslash}X@{}}
\toprule
Layer / object & Candidate id or relation & Law status & Role & Limit \\
\midrule
Lawful same-branch witness & \nolinkurl{intrinsic_naked_interface_witness_rewrite_effective_operator} & lawful & refined witness carrying the same-branch transport/readout law & not a public, cross-seed, or full-core carrier claim \\
Coarse public carrier & \nolinkurl{intrinsic_naked_interface_coarse} & signature\_only & public signature/readout shadow & does not carry the refined lawful transport structure \\
Witness-to-public relation & \nolinkurl{non_faithful_public_compression} & non-faithful & witness-to-public projection preserving public signature role & does not preserve the full law exactified on the witness \\
\bottomrule
\end{tabularx}
\end{table}

\subsection{Public-carrier upgrade attempts}

Three follow-up probes test whether the public layer can be upgraded, enriched, or
repaired into something stronger. None upgrades the coarse public candidate beyond
\texttt{signature\_only}.

The first probe is an authenticity-state backport, recorded as
\texttt{partially\_materialized}. It clarifies what authenticity-state information
reaches the active experiment surface and what remains missing, but it is diagnostic
context only: it does not rerun public-carrier law extraction and does not itself
upgrade the public carrier.

The second probe is a direct upgrade test. The enriched public-carrier row is rerun
against the same same-branch witness basis. The verdict is
\texttt{unchanged\_signature\_only}: the coarse public candidate remains
\texttt{signature\_only}; the primary witness remains lawful at the rewrite-effective
operator rung; and the compression relation remains
\texttt{non\_faithful\_public\_compression}. The enrichment neither moves the public
object into the lawful slot nor regresses the witness.

The third probe is a ceiling test. Orthogonal enrichment groups are combined into
increasingly strong public-facing surfaces, culminating in the
\texttt{full\_authenticity\_surface}. No admissible enrichment breaks the signature-only
ceiling, and no minimal successful public enrichment exists. The public shadow remains
coherent as a public signature surface, but the admissible enrichments tested here do
not turn it into the lawful object.

Had the public carrier been upgradeable into the lawful object by a modest authenticity
or enrichment correction, the witness/public distinction of the previous subsection
would have been a transient artifact of incomplete bookkeeping. The artifacts show the
opposite: the more carefully the public surface is audited, the more stable the
separation becomes. Spekkens's toy theory \citep{Spekkens2007} supplies a related
precedent in which operational statistics can arise from epistemic restrictions
without the public-facing description itself carrying the full underlying structure.

\subsection{What the split establishes}

Section~\ref{sec:live-surface-exactification} ended with the intrinsic branch as the
best live-surface carrier candidate with signature-level transport. This section
advances that result by identifying where law exactifies on the same branch: a lawful
same-branch witness exists and is the refined rewrite-effective-operator witness. The
coarse public carrier remains \texttt{signature\_only}, and the witness-to-public
relation is \texttt{non\_faithful\_public\_compression}; the lawful structure belongs
to the witness, not the public shadow.

Several questions remain open. Whether a stronger carrier surface survives only after
explicit protocol lifting and indexing, whether any cross-seed or full-core transfer
exists, and whether the split survives the later stress, honesty, and new-surface
panels are all deferred. The next section asks whether, on top of the lawful
same-branch witness, there is a further flagship-only protocol-indexed carrier surface
worth recording as a separate scoped positive.
 \section{Protocol-Indexed Flagship Carrier Surface}
\label{sec:protocol-indexed-flagship-carrier}

Section~\ref{sec:same-branch-law-and-public-shadow} fixed the two-level split: the
same-branch law lands on the refined witness
\nolinkurl{intrinsic_naked_interface_witness_rewrite_effective_operator}, while the
coarse public carrier remains a signature-level shadow under
\nolinkurl{non_faithful_public_compression}. The next question is whether the flagship
patch admits any further carrier-style strengthening. The strengthening lane answers
positively but with an explicit scope lock: it does not replace the witness/public-shadow
basis, and does not promote the intrinsic branch into a nonindexed carrier. It produces
a flagship-only surface whose surviving evidence remains protocol-coordinate indexed.

The framing is built into the cutover that opens the lane. Before the protocol lift is
constructed, the downstream basis is frozen as \texttt{witness\_site\_public\_shadow}
and the claim scope is restricted to the two flagship anchors. Later strengthening
steps are additions to that basis, not retroactive corrections of
Section~\ref{sec:same-branch-law-and-public-shadow}. The resulting surface is
\texttt{protocol\_indexed\_only}, and is stated as such.

\subsection{The autonomous protocol lift}

The first object in the strengthening lane is an autonomous lifted protocol surface
built on the same flagship patch as the lawful witness, with protocol coordinates made
first-class. The canonical artifact labels the result \texttt{indexed\_only}: the lift
is real and machine-readable, but survives only with protocol coordinates attached.

The lift is object-level, not a tagging trick. The registry records a finite lifted
witness object with explicit state, readout, and transport families: 1776 lifted states,
1776 readout entries, and six transport entries. It is stronger than the early
live-surface protocol-indexed branch of
Section~\ref{sec:live-surface-exactification}, which failed to instantiate coherently
on the direct core: there is now a finite object surface on which protocol
coordinates, not just public signatures, are part of the state description. The
treatment of protocol coordinates as first-class state data is methodologically close
to the geometric-phase and holonomy literature
\citep{Berry1984,Simon1983,AharonovAnandan1987}: the route or protocol traversed
becomes part of what distinguishes states on the surface.

The scope restriction matters. The lift covers only the two flagship stable ids. The
control-side stable ids
\nolinkurl{t27.control.null_baseline}, \nolinkurl{t27.control.memory_only}, and
\nolinkurl{t27.control.boundary_collapse} remain reference-only because the active
protocol/index artifacts do not provide exact source-backed bindings for them. The
paired artifact endpoints also remain explicit gap rows rather than being inferred
into the lifted object. The result is a flagship-patch construction with documented
coverage boundaries, not a broader direct-core or full-core surface.

The autonomous protocol lift is therefore not a nonindexed intrinsic carrier with a
few missing rows; it is a protocol-backed lift of the lawful same-branch witness on
the flagship patch.

\subsection{Future-only separation and observable predictive witnesses}

A loop-action probe, an observable-enrichment lane, and a predictive-witness
intersection sidecar together determine what evidence lives on the lifted flagship
patch. The loop-action probe finds \texttt{future\_only\_pairs} but \emph{not}
loop-action asymmetry: the summary artifact records
\texttt{loop\_action\_asymmetry\_status = not\_found}. The patch exhibits future-only
separation but does not realize the stronger Paper-2-style loop-action lane.

Future-only separation means that the lifted flagship surface carries nonempty pairs
that separate at the predictive future level while remaining collapsed at the current
quotient. That is a meaningful carrier-side signal: the protocol lift is not merely
relabeling the public shadow. Loop-action asymmetry would require a stronger
nontriviality pattern, which is absent. The flagship patch carries real future-separated
structure; the loop-action lane remains unavailable.

Observable enrichment salvages predictive-witness evidence across two families,
producing thirty predictive witnesses on the flagship patch. The lifted surface
therefore carries a concrete predictive-witness population and not merely abstract
future separation.

Control-side pairs remain source-limited or blocked. The screen-history overlay that
supports part of the provenance is treated as supporting evidence only; it is not
promoted to a separately protocol-audited candidate, and the later promotion gate
preserves that rule.

The predictive-witness intersection sidecar compares three populations: thirty
observable predictive pairs, seventy-eight lifted future-separated base pairs, and the
empty loop-action-asymmetry population. All thirty observable predictive pairs sit
inside the lifted future-separated population, and none sits inside the loop-action
population because no such population was found. The observable predictive witnesses
therefore track future separation rather than loop-action success. The positive signal
is future-only separation plus observable predictive-witness salvage on a lifted
flagship patch, with no loop-action asymmetry.

\subsection{Honesty on the lifted flagship patch}

Once the lifted object and its predictive population are in place, a lifted honesty rerun
and a joint attack matrix ask whether the same flagship patch survives the honesty fronts
needed for a serious scoped carrier claim. The answer is again positive and again
carefully bounded. The lifted section/protocol gate returns
\texttt{lifted\_honesty\_status = flagship\_protocol\_robust}. The audited candidate is
\nolinkurl{intrinsic_naked_interface_witness_autonomous_protocol_lift}; the audited
stable ids are exactly the two flagships; the section verdict is \texttt{section\_robust};
and the protocol verdict is \texttt{protocol\_robust}.

The honesty results show that the lifted flagship surface is not a fragile reindexing
trick: the object that carries future-separated predictive witnesses also survives the
section and protocol reruns on the audited patch. The inherited gap rows consist of
the three direct-core controls and the two paired artifact endpoints, outside the
audit scope.

The joint attack matrix returns \texttt{both\_fronts\_closed}: the predictive front is
closed, the protocol front is closed, and the loop-action front remains
\texttt{not\_found}. The scope is \texttt{flagship\_only}, and the recommended next
move is the intrinsic promotion gate. The missing loop-action lane and the unresolved
control scope keep the result from being advertised as broader than flagship.

The control scope contract is still unresolved at this stage; later sections test what
happens when controls and nearby positives are pushed harder. The flagship patch is
coherent enough to survive both supported fronts while broader transfer questions are
deferred.

\subsection{The intrinsic promotion gate}

The intrinsic promotion gate decides what kind of survivor the strengthening lane
yields. The audited surviving candidate is
\nolinkurl{intrinsic_naked_interface_witness_autonomous_protocol_lift}, and its final
promotion status is \texttt{protocol\_indexed\_only}.

Two subordinate audits explain why the verdict stops there. The second-stage audit
returns \texttt{identity\_only}: only the identity proposal survives, and every
nonidentity merge proposal collapses deterministic transport or witness retention, so
the lane does not earn a deeper promoted object by merging away protocol distinctions.
The signature-update audit returns \texttt{indexed\_signature\_only}, with the
strongest blocking reason being that no overlap witness remains separated once
coordinate tokens are stripped. The positive evidence disappears when protocol-coordinate
indexing is removed.

The lifted object is real, the predictive and protocol fronts close on the flagship
patch, and the surviving evidence depends on protocol coordinates. The verdict is
therefore stronger than ``nothing survives'' and weaker than ``the intrinsic carrier has
promoted outright'': a scoped, protocol-indexed survivor has been earned.

The gate is explicit about the operational basis. \texttt{paper\_may\_upgrade\_basis}
is \texttt{false}, and the recommended basis remains
\texttt{witness\_site\_public\_shadow}. The witness/public-shadow pair remains the
basis for same-branch law and public readout; the autonomous protocol lift is recorded
as an additional flagship-only surface that survives only with indexing attached.

The active screen-history overlay is kept at \texttt{supporting\_evidence\_only}: it
explains part of the predictive-witness provenance on the flagship patch but is not a
separately protocol-audited promotion candidate. There is one audited surviving
candidate, with status \texttt{protocol\_indexed\_only}. The analogy to operational
contextuality is therefore only methodological \citep{Spekkens2005}: once the
relevant context is made explicit, the surviving description here also has to retain
that context---the protocol coordinates---as part of the object.

Table~\ref{tab:protocol-indexed-strengthening} summarizes the strengthening lane step by
step.

\begin{table}[t]
\centering
\footnotesize
\caption{Flagship strengthening lane after the public-shadow cutover. The lane closes
predictive and protocol fronts on the flagship patch, but the intrinsic promotion gate
leaves the survivor protocol-indexed rather than nonindexed.}
\label{tab:protocol-indexed-strengthening}
\begin{tabularx}{\linewidth}{@{}>{\raggedright\arraybackslash}p{0.14\linewidth}>{\raggedright\arraybackslash}p{0.20\linewidth}>{\raggedright\arraybackslash}p{0.22\linewidth}>{\raggedright\arraybackslash}p{0.20\linewidth}>{\raggedright\arraybackslash}X@{}}
\toprule
Step & Surface tested & Positive finding & Limitation & Result used in this paper \\
\midrule
Witness/public-shadow cutover & witness/public-shadow basis & downstream basis fixed as \nolinkurl{witness_site_public_shadow}; claim scope restricted to flagships & no public-side law cutover and no broader basis replacement & starting basis for the strengthening lane \\
Autonomous protocol lift & protocol-backed flagship patch & finite lifted object built; \texttt{indexed\_only} & controls and paired endpoints remain reference-only gaps & real object-level flagship surface \\
Future-only probe & lifted flagship patch & \texttt{future\_only\_pairs} found & loop-action asymmetry not found & future-separated surface, not a loop-action lane \\
Observable enrichment & flagship predictive population & 30 predictive witnesses salvaged across 2 families & control-side rows remain source-limited; overlay is provenance only & observable predictive support on the flagship patch \\
Predictive-witness intersection & future-separated flagship pairs & all 30 observable predictive pairs lie inside the 78 future-separated base pairs & none lie in a loop-action-asymmetry population & predictive witnesses track future separation, not loop-action success \\
Lifted honesty rerun & audited lifted candidate on the two flagships & section robust and protocol robust & audit remains flagship-only with inherited control/endpoint gaps & flagship protocol robust \\
Joint attack matrix & same flagship patch & predictive and protocol fronts both closed & loop-action not found; control scope unresolved & scoped closure before promotion \\
Intrinsic promotion gate & audited autonomous lift & survivor recorded as \texttt{protocol\_indexed\_only} & identity-only second stage; indexed signature only; no basis upgrade & scoped positive carrier surface used later in the paper \\
\bottomrule
\end{tabularx}
\end{table}

\subsection{What the protocol-indexed survivor establishes}

Section~\ref{sec:same-branch-law-and-public-shadow} established the lawful same-branch
witness and the public signature shadow. The present section adds a separate
flagship-only carrier-style surface whose positive evidence survives only with
protocol coordinates attached.

The positive findings on that patch are: a lifted flagship object, future-only
separation, observable predictive witnesses, robust section and protocol honesty,
closure of both supported fronts in the joint attack matrix, and one surviving
audited candidate at the promotion gate.

The limitations are structural. The survivor is \texttt{flagship\_only} and
\texttt{protocol\_indexed\_only}; the operational basis remains
\texttt{witness\_site\_public\_shadow}; no nonindexed intrinsic carrier, broader
direct-core or cross-seed transfer, or lawful public carrier has been established; and
future-only separation is not converted into loop-action asymmetry.

Once a flagship-only, protocol-indexed survivor has been identified, the next question
is whether any of its positivity transfers beyond the flagship patch or survives the
scope contracts applied to nearby positives and controls.
 \section{Cross-Seed Reach and Scope Contracts}
\label{sec:cross-seed-reach-and-scope-contracts}

Section~\ref{sec:protocol-indexed-flagship-carrier} established a tightly scoped
positive surface: a lawful same-branch witness, a public signature shadow, and a
flagship-only protocol-indexed strengthening lane. Whether that surface transfers
beyond the anchors is addressed in two steps. A cross-seed canonicalization pass
measures transfer on the frozen non-flagship panel; a scope-contract pass then resolves
what kind of non-flagship evidence each panel represents. The combined result records
meaningful support alignment beyond the anchors without upgrading that alignment into
exact cross-seed readout or law transfer.

\subsection{Cross-seed canonicalization}

The cross-seed reach test returns \texttt{canonicalization\_status=partial} on the
canonical post-strengthening basis. The question is asked relative to the living basis
of the preceding sections---the lawful same-branch witness, the public signature
shadow, and the flagship-only protocol-indexed strengthening lane, with the downstream
basis still fixed as \texttt{witness\_site\_public\_shadow}---not relative to a
hypothetical stronger carrier.

On that basis, the two flagship anchors remain the only seeds carrying exact readout and
law transfer. The four nearby positives do not disappear from the corpus. On the
contrary, they remain visible enough to register as \texttt{support\_only\_alignment}.
That is the right phrase to emphasize here. The nearby-positive panel is not empty, and
it is not reclassified as a control panel. It preserves meaningful alignment with the
frozen seed/reduction layer. But that alignment does not materialize on the canonical
lawful witness or on the coarse public carrier as exact readout or transport-law
transfer.

All four nearby positives are classified as \texttt{support\_only\_alignment}. None
transfers to the primary lawful witness, none to the coarse public carrier; no exact
readout canonicalization and no transport-law transfer is recorded. The result is
therefore partial canonicalization: more than noise beyond the anchors, less than
carrier transfer.

The inherited strengthening-lane limitations remain active. The strengthening lane is
\texttt{flagship\_only} and \texttt{protocol\_indexed\_only} and serves as comparison
context rather than an upgraded full-core basis. The control scope contract is still
unresolved at this stage, and the loop-action disposition remains
\texttt{not\_found\_on\_current\_surface}. These limits prevent support alignment from
being used as if it already carried the same status as the anchors.

The public side records zero public-carrier transfer beyond the anchors. This is
consistent with Section~\ref{sec:same-branch-law-and-public-shadow}: the public object
remains a signature-level shadow, and nothing in the cross-seed panel upgrades it into
a broader law-bearing surface. The lawful witness is real and intact on the same
branch but is not transferred to the nearby-positive panel.

\subsection{Non-flagship scope contracts}

The scope-contract pass returns
\texttt{scope\_contract\_status=resolved\_with\_real\_limitations}. It assigns every
non-flagship panel a clear status, separating three reasons why a seed or panel should
not be promoted into a broader carrier claim.

\texttt{null\_by\_design}. All four control-side cases are signed off this way: the
null-baseline control, the memory-only control, the paired artifact control, and the
boundary-collapse control. They are controls, not positive-transfer failures. The
paired artifact control remains a paired logical panel excluded from ordinary transfer
counts.

\texttt{real\_scope\_limitation}. This applies to the four nearby positives. They are
not null controls and not administrative artifacts: they are evidence of the current
limit of transfer, aligning at support level but not landing on the lawful witness or
on the public signature shadow with exact readout or law. Full-core transfer therefore
remains unlicensed after the scope contract, even though the positive flagship
surfaces remain intact.

\texttt{positive\_transfer}. Only the support-only erasure-recovery anchor receives
this label, and only as \emph{support-only transfer}: the row remains source-backed
and useful but does not upgrade to exact readout transfer or transport-law transfer.
This is why the support-only panel is kept separate from both the nearby positives and
the controls.

Table~\ref{tab:scope-contracts} summarizes the cross-seed and scope-contract classification
of the non-flagship panel.

\begin{table}[t]
\centering
\footnotesize
\caption{Cross-seed and non-flagship scope contracts. The nearby positives align only at
support level and remain real scope limitations; controls are signed off as
null-by-design, and the support-only row remains support-only.}
\label{tab:scope-contracts}
\begin{tabularx}{\linewidth}{@{}>{\raggedright\arraybackslash}p{0.19\linewidth}>{\raggedright\arraybackslash}p{0.18\linewidth}>{\raggedright\arraybackslash}p{0.17\linewidth}>{\raggedright\arraybackslash}p{0.16\linewidth}>{\raggedright\arraybackslash}X@{}}
\toprule
Panel / role & Count or stable ids & Cross-seed classification & Scope contract & Interpretation for carrier claim \\
\midrule
Flagship anchors & 2 anchors & exact readout and law transfer & anchor basis, not part of non-flagship contract counts & only positive source panel for the current carrier claim \\
Nearby positives & 4 nearby positives & support-only alignment & real scope limitation & meaningful alignment beyond anchors, but no exact cross-seed readout or law transfer \\
Null-baseline control & \nolinkurl{t27.control.null_baseline} & exact basis materialization as a control row & null by design & control consistency, not positive transfer \\
Memory-only control & \nolinkurl{t27.control.memory_only} & exact basis materialization as a control row & null by design & control consistency, not positive transfer \\
Artifact control & \nolinkurl{t27.control.artifact} & support-only alignment on the paired logical panel & null by design & paired panel kept separate from ordinary transfer counts \\
Boundary-collapse control & \nolinkurl{t27.control.boundary_collapse} & exact basis materialization as a control row & null by design & boundary control, not positive transfer \\
Support-only erasure-recovery anchor & \nolinkurl{t27.support.erasure_recovery_anchor} & support-only alignment & positive transfer & positive only as support-only transfer; not exact readout or law transfer \\
\bottomrule
\end{tabularx}
\end{table}

The scope-contract vocabulary is more informative than simple pass/fail language.
Controls are not demoted from positive transfer because they were never meant to be
positive transfer cases. The nearby positives are not nullified, because their support
alignment is real. And the support-only row is not overpromoted merely because it
survives as a useful positive transfer in its restricted sense. The result is
positive on the flagship basis, aligned beyond the anchors at support level, and
blocked from full-core transfer by the nearby-positive panel.

\begin{figure}[t]
\centering
\resizebox{0.90\linewidth}{!}{\begin{tikzpicture}[
  font=\small,
  node distance=0.9cm and 0.9cm,
  group/.style={
    draw,
    rounded corners,
    align=center,
    fill=black!4,
    text width=0.28\linewidth,
    minimum height=1.3cm,
    inner sep=5pt
  },
  sidegroup/.style={
    draw,
    rounded corners,
    align=center,
    fill=black!2,
    text width=0.28\linewidth,
    minimum height=1.3cm,
    inner sep=5pt
  }
]
\node[group, text width=0.34\linewidth] (flagship) {Flagship anchors (2)\\scoped protocol-indexed positive carrier surface};
\node[group, below left=of flagship] (nearby) {Nearby positives (4)\\support-only alignment\\real scope limitations};
\node[group, below=of flagship] (controls) {Ordinary controls (3)\\null-by-design};
\node[group, below right=of flagship] (support) {Support-only row (1)\\support-only positive transfer};
\node[sidegroup, below=of controls] (artifact) {Paired artifact panel\\separate logical control\\not ordinary support};
\end{tikzpicture}
 }
\caption{Scope panels after the scope-contract pass. Only the two flagship anchors
support the scoped protocol-indexed carrier surface; the nearby positives remain
support-only alignment and real scope limitations, the ordinary controls remain
null-by-design, the support-only row remains support-only, and the paired artifact panel
stays separate from ordinary support counts.}
\label{fig:scope-panel}
\end{figure}

Figure~\ref{fig:scope-panel} provides the same panel split in compact visual form and
keeps the paired artifact control outside the ordinary seed interpretation used for the
flagship, nearby, control, and support-only rows.

\subsection{What remains licensed after the cross-seed test}

What remains licensed is the basis already established by
Sections~\ref{sec:same-branch-law-and-public-shadow} and
\ref{sec:protocol-indexed-flagship-carrier}: the lawful same-branch witness, the public
signature shadow under non-faithful compression, and the flagship-only protocol-indexed
strengthening lane. Broader cross-seed or full-core transfer is not licensed; the
nearby-positive panel blocks it as a \emph{real scope limitation} rather than by
collapsing into noise.

The run reaches beyond the anchors only as support alignment; the stronger carrier
content remains tied to the flagship basis. The public side does not inherit law on
the nearby-positive panel, and the protocol-indexed flagship surface does not become a
broader nonindexed transfer because some support alignment exists outside the anchors.

The next section asks whether the scoped basis survives the falsifier panels:
denominator stress, surface expansion, refreshed honesty, and adjacent-surface search.
 \section{Stress, Expansion, Honesty, and New-Surface Panels}
\label{sec:stress-expansion-honesty-new-surface}

Sections~\ref{sec:same-branch-law-and-public-shadow}--\ref{sec:cross-seed-reach-and-scope-contracts}
fix the carrier result at a particular level. Four falsifier panels then ask whether
that bounded result survives the standard stress tests: collapse of the licensed
witness, fragility of the witness/public-shadow composite under stronger overclaim
demands, or failure of broader globalized or adjacent-surface claims. Only the last
two occur. The methodological commitment is the standard falsifiability discipline
\citep{Popper1959}: each panel states in advance what would discredit the licensed
claim.

\subsection{Denominator stress}

The denominator-stress panel applies strict denominator tests to the canonical
post-strengthening basis and returns a fragile denominator verdict. What is fragile is
not the licensed lawful witness: the primary lawful witness is classified
\emph{robust}. The licensed law surface survives; the observed failures are confined
to coarse-public weakness, non-faithful compression, support-only overpromotion, and
full-core overclaim.

The coarse public carrier is \texttt{fragile}: demanding stronger denominator behavior
only re-exposes its signature-level status. The witness/public-shadow composite is
also \texttt{fragile}, not because it ceases to function on its licensed basis, but
because it depends on non-faithful compression and does not generalize into a broader
full-core law.

The support-only row is a third kind of fragile object: it survives only as
support-only transfer and becomes fragile as soon as stricter exact-readout or law
semantics are imposed. This confirms the scope-contract classification---the row is
usable but not an exact or lawful transfer row.

The full-core challenge collapses, not because the witness breaks but because the
overclaim sentinel breaks where the scope contract says it should: at the
nearby-positive real scope limitations and at the refusal to promote support-only
transfer into law transfer. The licensed witness survives strict denominator pressure;
the broader claim surfaces do not.

\subsection{Surface expansion stress}

The surface-expansion panel expands the canonical basis through the repo-native
families closest to same-branch enrichment, public authenticity enrichment, lifted
flagship protocol indexing, and route filtering. Its branch decision is
\texttt{expansion\_status=fragile}; the question is how that fragility is
distributed.

The primary lawful witness remains \emph{robust}: the exercised same-branch
shell-compatible expansions do not dislodge it. The lawful witness remains the stable
center of the post-law-audit basis.

The coarse public carrier remains fragile under expansion, as expected after
Section~\ref{sec:same-branch-law-and-public-shadow}: public-authenticity and
non-same-branch enrichment attempts continue to expose the public side as weaker than
the witness side. The witness/public-shadow composite is also fragile, because
expansion reopens the same compression and public-side weakness without turning the
composite into a broader law-bearing carrier.

The support-only row is recorded as \emph{robust} under the exercised expansion
families, but only as support-only transfer. Robustness at that level does not upgrade
it into exact readout or lawful transport transfer.

The strengthening lane is recorded \texttt{fragile}. It remains comparison support
only---\texttt{flagship\_only}, \texttt{protocol\_indexed\_only}---not an upgrade into
a same-branch replacement basis.

The route-filtering family stays unavailable. The panel keeps three kinds of expansion
outcome separate: same-branch shell-compatible families that the witness survives,
branch-break or new-surface families that do not count as same-branch witness failure,
and route-unavailable probes recorded as unavailable rather than treated as evidence
either way.

\subsection{Section and protocol honesty}

The refreshed honesty gate revisits the section-and-protocol audit after the earlier
panels: the refresh mode is \texttt{refreshed\_from\_existing}, and the status is
unchanged. The branch decision remains
\texttt{honesty\_status=site\_not\_probe\_fails\_protocol\_honesty}. The failure
attaches to the global/full site-not-probe claim, not to the licensed scoped witness
basis.

The target-level map: the primary lawful witness is \texttt{protocol\_robust}; the
support-only positive transfer row is \texttt{protocol\_robust} in its narrow role; the
flagship protocol-indexed strengthening lane remains \texttt{protocol\_indexed\_only};
the coarse public carrier and the witness/public-shadow composite are
\texttt{section\_fragile}; and the full-core claim challenge is
\texttt{artifact\_collapse}.

The gate summary: the licensed witness/public-shadow basis stays usable, but the
global/full site-not-probe claim fails protocol honesty because the nearby positives
remain real scope limitations and the strongest lifted lane remains
protocol-indexed-only. The panel is not invalidating the scoped basis; it is refusing
to let the scoped basis be rewritten as a broader full-core or global claim.

The paired panel remains \texttt{separate\_explicit}, so it is not smuggled into
ordinary panel counts. Route filtering remains unavailable, so the honesty refresh has
no license to repair the full/global claim through a probe that the active surfaces
cannot evaluate. Both facts contribute to the \texttt{artifact\_collapse} protocol
verdict on the broader overclaim. The discipline parallels the loophole-closure
regime in experimental quantum foundations---for example, the loophole-free
Bell-test programme \citep{Hensen2015}---and, more loosely, the no-go literature
that constrains which epistemic readings can reproduce the target statistics under
additional assumptions \citep{PuseyBarrettRudolph2012}.

The licensed basis survives honesty, the primary lawful witness survives honesty, and
the coarse public carrier remains a signature-level public side. The honesty gate is
not a destructive reset; it shows that the scoped basis is stable when it is not
overdescribed as a broader site-not-probe or full-core claim.

\subsection{Adjacent-surface search}

The adjacent-surface search asks whether the limitations exposed by the preceding
panels can be repaired by a small adjacent-surface escalation. Its branch decision is
\nolinkurl{new_surface_status=no_exactification}.

The search selects three repo-native candidates closest to the current basis: a
selection-support-enriched same-branch surface, a public authenticity-enriched
surface, and the lifted flagship protocol-indexed surface carried forward from the
strengthening lane and the expansion panel. It skips the observable predictive-witness
overlay because that overlay remains supporting provenance rather than an independent
exactification surface, and records the route-filtering probe as unavailable rather
than silently absent. The controlled-extension idiom is methodologically reminiscent
of scope-restricted effective theories \citep{Crowther2016} and of operational
reconstructions that enlarge the state description through purification
\citep{ChiribellaDArianoPerinotti2010}.

None of the selected adjacent surfaces materially improves cross-seed or broader
full-core transfer; none upgrades the coarse public side; the lawful witness is
preserved as the comparison baseline throughout. The best selected surface is the
lifted flagship protocol-indexed surface, but even there the status is
\texttt{insufficient\_signal}.

The selection-support-enriched same-branch surface is shell-compatible and
same-branch but does not repair the nearby-positive limitation; the
public-authenticity-enriched surface is the closest public-side probe but does not
improve the public shadow into a stronger carrier; and the lifted flagship
protocol-indexed surface remains tied to a flagship-only indexed scope. The skipped
overlay remains provenance only, and route filtering remains unavailable.

Table~\ref{tab:falsifier-panels} summarizes the four falsifier panels.

\begin{table}[t]
\centering
\footnotesize
\caption{Stress, honesty, and adjacent-surface summary. The licensed witness survives the
falsifier panels, while broader full-core, public-side, and adjacent-surface upgrades do
not land.}
\label{tab:falsifier-panels}
\begin{tabularx}{\linewidth}{@{}>{\raggedright\arraybackslash}p{0.14\linewidth}>{\raggedright\arraybackslash}p{0.18\linewidth}>{\raggedright\arraybackslash}p{0.17\linewidth}>{\raggedright\arraybackslash}p{0.16\linewidth}>{\raggedright\arraybackslash}X@{}}
\toprule
Panel & Branch decision & Licensed-basis result & Overclaim / new-surface result & Interpretation used in this paper \\
\midrule
Denominator stress & \texttt{fragile} & primary lawful witness robust; public shadow and compression composite fragile & full-core challenge collapses; support-only overpromotion stays fragile & strict denominators do not erase the licensed witness, but they expose broader overclaim weakness \\
Surface expansion & \texttt{fragile} & lawful witness robust on same-branch shell-compatible expansions; support-only row remains support-only & public-side and composite fragility persist; route filtering unavailable; full-core challenge collapses & expansion preserves the scoped basis without globalizing it \\
Section/protocol honesty & site-not-probe fails protocol honesty & licensed witness/public-shadow basis stays usable; primary witness survives honesty & full/global site-not-probe claim fails; paired panel stays separate; route filtering unavailable & honesty failure is attached to the broader claim, not to the scoped licensed surfaces \\
Adjacent-surface search & \texttt{no\_exactification} & lawful witness preserved as baseline; best adjacent surface still insufficient & no cross-seed/full-core improvement; no public-side upgrade; no exactified new surface & repo-native nearby surfaces do not repair the earlier scope limits \\
\bottomrule
\end{tabularx}
\end{table}

\subsection{Synthesis of the falsifier panels}

The four panels identify the honest boundary of the positive result. The lawful
same-branch witness remains standing. The public side remains a signature shadow
rather than a stronger law-bearing carrier. The flagship-only protocol-indexed
strengthening lane remains real but scoped. What fails are the attempts to reinterpret
those surfaces as a full-core or wider-transfer result, and the adjacent-surface
attempts to bypass the nearby-positive limitations.

The basis survives the stress, expansion, honesty, and adjacent-surface tests needed
to count as a serious scoped result, while the broader claims fail in reproducible and
interpretable ways. Scope restriction of this kind is, in the philosophy-of-physics
tradition, the normal consequence of singular limits and effective-theory boundaries
\citep{Batterman2002}. The next section asks whether any downstream mathematical
structure persists on the same scoped carrier surface.
 \section{Scoped Quadratic Carrying and Proto-Event Tests}
\label{sec:scoped-quadratic-carrying-and-proto-event-tests}
\label{sec:quadratic-carrying-and-proto-event-pilot}

Sections~\ref{sec:cross-seed-reach-and-scope-contracts} and
\ref{sec:stress-expansion-honesty-new-surface} leave a sharply bounded result: the
lawful same-branch witness survives, the coarse public carrier remains signature-only,
the flagship strengthening lane remains protocol-indexed-only, and no cross-seed,
full-core, or adjacent-surface upgrade lands. The scoped quadratic and proto-event
probes ask a narrower downstream question: does any secondary structure continue to
live on the earned scoped carrier? The answer is yes, at the same flagship,
protocol-indexed scope.

\subsection{Opening the quadratic branch at scoped carrier scope}

The quadratic branch opens under the run-gate policy
\nolinkurl{earned_positive_carrier_any_survived_class}: the branch does not wait for
cross-seed or full-core exactification but asks whether any honestly survived carrier
surface exists on which quadratic structure can be tested without misstating the
claim. The run condition is met because an earned positive carrier survives in scoped
form.

The surviving carrier surface used by the quadratic branch is
\nolinkurl{intrinsic_naked_interface_witness_autonomous_protocol_lift}, with survival
class \texttt{protocol\_indexed\_carrier\_only} and scope class \texttt{flagship\_only}.
The quadratic scope is therefore \texttt{flagship\_only\_protocol\_indexed}: the same
scoped survivor from Section~\ref{sec:protocol-indexed-flagship-carrier}, not a new
basis.

The quadratic branch does not replace the witness/public-shadow basis. The primary
lawful witness remains
\nolinkurl{intrinsic_naked_interface_witness_rewrite_effective_operator}; the coarse
public carrier remains \nolinkurl{intrinsic_naked_interface_coarse} with
\texttt{signature\_only} status; and the relation between them remains
\texttt{non\_faithful\_public\_compression}. The branch opens because a scoped carrier
has survived on top of that basis, not because the basis itself has been widened.

A positive quadratic result here does not reverse the earlier cross-seed, scope,
stress, honesty, or new-surface limitations. It says that the earned scoped carrier is
substantial enough to carry some downstream structure.

\subsection{Claim-bearing quadratic families}

With that gate opened, the branch returns \texttt{quadratic\_status=carried}. The carried
families are three claim-bearing quadratics:
\begin{enumerate}[label=\arabic*.,leftmargin=2.0em,itemsep=0.2em]
\item \texttt{q\_readout\_l2\_self},
\item \texttt{q\_readout\_gap},
\item \texttt{q\_coarsening\_consistency}.
\end{enumerate}
Each is recorded as \texttt{carried} on the same scoped carrier surface. The three
families are the claim-bearing quadratic probes that factor through the surviving
flagship-only, protocol-indexed carrier and remain compatible with the scoped honesty
and stress panels.

The carried quadratics do not float independently of the witness structure. They are
carried because the primary lawful witness remains intact, because the autonomous
lifted surface remains section- and protocol-robust at its own scope, and because the
promotion audit leaves the surviving row identity-stable but still indexed. The branch
inherits the carrier result's scope and lives there.

The bilinear family \nolinkurl{b_aligned_readout_inner_product} remains
\texttt{comparison\_only}. The branch opens the comparison-side sub-branch, but the
bridge summary still classifies that row as scoped comparison only. The bilinear family
is retained as comparison-only and is not promoted into a broader bridge claim.

No cross-seed quadratic carrying is claimed, and no broader bridge claim is claimed
from the bilinear family. The quadratics carry because an earned positive carrier
survives in scoped form, not because a stronger carrier surface was discovered.
Gleason's theorem \citep{Gleason1957} is a canonical example of measure-theoretic
structure derived under non-contextual assignment assumptions; our scoped result is
compatible with but much weaker than that kind of derivation.

\subsection{Proto-event/test pilot}

The proto-event/test pilot asks whether any proto-event/test families survive on the
same scoped carrier. The relevant operational precedent is the effect-algebra
formalism \citep{FoulisBennett1994}, which treats scoped test-and-event structures
without requiring a full event algebra. The branch decision is
\texttt{proto\_event\_status=present}. The carrier surface remains
\nolinkurl{intrinsic_naked_interface_witness_autonomous_protocol_lift}. Its survival
class remains protocol-indexed carrier only, its scope class remains flagship-only,
and the pilot scope remains flagship-only and protocol-indexed.

The pilot tests three candidate proto-event/test families: two are present, one is
absent, none merely not evaluable. The present families are
\nolinkurl{flagship_quadratic_levelset_event}, a levelset-style test supported by the
carried quadratics, and \nolinkurl{flagship_transition_stability_test}, a
transition-stability test that remains present because the surface is still section-
and protocol-robust even though it is only identity-stable and does not promote to a
broader carrier class. The absent family,
\nolinkurl{flagship_protocol_indexed_readout_agreement_event}, does not survive
because the readout relation remains too asymmetric across the two flagship rows to
support a surviving readout-agreement event family.

Protocol indexing remains required. The paired panel is \texttt{not\_touched}. Nearby
positives, controls, and the support-only row are not treated as ordinary positive
proto-event support: the pilot is scoped to the surviving flagship-only,
protocol-indexed carrier, and the nearby positives remain real scope limitations
rather than positive event rows.

Table~\ref{tab:scoped-quadratic-protoevent} summarizes the quadratic and proto-event
outcomes.

\begin{table}[t]
\centering
\footnotesize
\caption{Scoped quadratic carrying and proto-event/test pilot. All positive entries remain
at the flagship protocol-indexed carrier scope; no broader bridge or general
event-geometry claim is made.}
\label{tab:scoped-quadratic-protoevent}
\begin{tabularx}{\linewidth}{@{}>{\raggedright\arraybackslash}p{0.22\linewidth}>{\raggedright\arraybackslash}p{0.15\linewidth}>{\raggedright\arraybackslash}p{0.19\linewidth}>{\raggedright\arraybackslash}p{0.18\linewidth}>{\raggedright\arraybackslash}X@{}}
\toprule
Item & Status & Scope & Role & Non-claim / limitation \\
\midrule
Carried carrier surface & survived scoped carrier & flagship-only, protocol-indexed & earned protocol-indexed carrier basis for downstream probes & not a cross-seed or full-core carrier surface \\
\nolinkurl{q_readout_l2_self} & carried & flagship-only, protocol-indexed & claim-bearing quadratic family & remains scoped to the survived carrier \\
\nolinkurl{q_readout_gap} & carried & flagship-only, protocol-indexed & claim-bearing quadratic family & does not widen the carrier claim \\
\nolinkurl{q_coarsening_consistency} & carried & flagship-only, protocol-indexed & claim-bearing quadratic family & carried only on the indexed survivor \\
\nolinkurl{b_aligned_readout_inner_product} & comparison-only & flagship-only, protocol-indexed & bilinear comparison family & not promoted into a broader bridge claim \\
Quadratic levelset event & present & flagship-only, protocol-indexed & scoped proto-event/test family & not a general event family \\
Transition stability test & present & flagship-only, protocol-indexed & scoped proto-event/test family & depends on the same scoped carrier robustness \\
Readout-agreement event & absent & flagship-only, protocol-indexed & attempted readout-agreement proto-event & no surviving exact agreement event on the current scoped carrier \\
\bottomrule
\end{tabularx}
\end{table}

\subsection{Scope of the secondary structure}

Neither the quadratic branch nor the proto-event pilot creates a new top-level
verdict. They add secondary structure on the same scoped carrier. Every positive entry
in the section remains \texttt{flagship\_only} and \texttt{protocol\_indexed\_only}:
the carried quadratic families remain scoped, the two present proto-event/test
families remain scoped, the comparison-only bilinear family remains comparison-only,
and the absent readout-agreement row remains absent.

The results add evidence that the surviving scoped carrier is mathematically
nontrivial without widening the carrier claim, turning the bilinear comparison into a
bridge, or turning proto-events into a general theory. The next section treats them as
secondary flags supporting the final atlas: scoped quadratic carrying and scoped
proto-events present.
 \section{Final Gate and Carrier Atlas}
\label{sec:final-gate-and-carrier-atlas}
\label{sec:final-gate}

The previous sections separated surfaces that a looser narrative might have merged:
the intrinsic live-surface candidate of Section~\ref{sec:live-surface-exactification},
the lawful same-branch witness and public signature shadow of
Section~\ref{sec:same-branch-law-and-public-shadow}, the flagship-only
protocol-indexed survivor of Section~\ref{sec:protocol-indexed-flagship-carrier}, the
cross-seed and honesty limitations of
Sections~\ref{sec:cross-seed-reach-and-scope-contracts} and
\ref{sec:stress-expansion-honesty-new-surface}, and the scoped quadratic and
proto-event diagnostics of
Section~\ref{sec:quadratic-carrying-and-proto-event-pilot}. The final gate aggregates
those threads into a single carrier verdict and states which claims are licensed.

The aggregation is sharper than a yes/no gate. The run does not end with full carrier
exactification nor with the absence of any carrier surface; the gate records a
stronger negative and a weaker positive together. The strongest public, cross-seed,
and full-core carrier claim does not land; a flagship-only, protocol-indexed carrier
surface survives on top of the lawful same-branch witness. This section states that
outcome and shows how the carrier atlas, the C1--C7 criteria, and the claim-permission
map fit together.

\subsection{Final gate verdict}

The final gate verdict is \texttt{protocol\_indexed\_carrier\_only}. The strongest
claim verdict is \texttt{failed}, and the survived positive surface verdict is
\texttt{present}. These three lines must be read together: the strongest-claim row
alone would make the run sound purely negative, while the survived-positive row alone
would invite an overreading into a stronger carrier verdict than the artifacts
support.

The surviving positive carrier surface is
\nolinkurl{intrinsic_naked_interface_witness_autonomous_protocol_lift}, with scope
\texttt{flagship\_only} and survival class \texttt{protocol\_indexed\_carrier\_only}.
The scoped survivor does not replace the same-branch witness of
Section~\ref{sec:same-branch-law-and-public-shadow}: the primary lawful witness
remains \nolinkurl{intrinsic_naked_interface_witness_rewrite_effective_operator}, and
the coarse public carrier remains \nolinkurl{intrinsic_naked_interface_coarse} with
\texttt{signature\_only} status under \texttt{non\_faithful\_public\_compression}. The
central split is preserved, not flattened.

No strongest new carrier exactifies at public, cross-seed, or full-core scope. The run
does preserve a carrier surface stronger than a mere descriptor or bookkeeping
trace---a flagship-only, protocol-indexed survivor that remains honest about its
indexing and scope. The lawful same-branch witness is present, the public signature
shadow is retained in its coarse role, and the downstream quadratic and proto-event
rows are available only on the scoped survivor.

The strongest claim is blocked by explicit factors in the gate record: cross-seed
canonicalization remains partial; the scope contract keeps the four nearby positives
as real scope limitations; the refreshed honesty gate preserves the collapse of the
global/full site-not-probe claim; the coarse public carrier remains signature-only;
the adjacent-surface search earns no upgrade; and the C3 audit reduces the richest
surviving scoped carrier back to an already-landed simpler lane for final-gate
purposes.

Table~\ref{tab:final-gate-summary} makes that split explicit by keeping the strongest
claim verdict and the survived positive surface verdict visible in the same summary.

\begin{table}[t]
\centering
\footnotesize
\caption{Final gate summary. The strongest public/full-core/cross-seed carrier claim
fails, but a flagship-only protocol-indexed carrier surface survives together with a
lawful same-branch witness.}
\label{tab:final-gate-summary}
\begin{tabularx}{\linewidth}{@{}>{\raggedright\arraybackslash}p{0.24\linewidth}>{\raggedright\arraybackslash}p{0.19\linewidth}>{\raggedright\arraybackslash}p{0.19\linewidth}>{\raggedright\arraybackslash}X@{}}
\toprule
Item & Status & Scope & Consequence \\
\midrule
Final gate verdict & \nolinkurl{protocol_indexed_carrier_only} & flagship-only, protocol-indexed & strongest surviving carrier surface is positive but explicitly scoped \\
Strongest carrier claim & failed & public/full-core/cross-seed & no strongest new-carrier verdict is licensed \\
Survived positive surface & present & flagship-only, protocol-indexed & a weaker but real carrier surface remains in the atlas \\
Coarse public carrier & \texttt{signature\_only} & public shadow & stable public signature surface, not carrier law \\
Lawful same-branch witness & lawful & same branch & primary lawful survivor of the run \\
Scoped protocol-indexed carrier & survived scoped carrier & flagship-only, protocol-indexed & strongest positive carrier row retained by the gate \\
Cross-seed test & partial / not licensed & nearby-positive panel & support alignment only; no exact readout or law transfer \\
New-surface search & no exactification & adjacent-surface search & no alternate surface upgrades the verdict \\
Quadratic carrying & scoped carried & flagship-only, protocol-indexed & secondary positive flag on the survived carrier \\
Proto-event/test pilot & scoped present & flagship-only, protocol-indexed & secondary positive flag on the same survived carrier \\
\bottomrule
\end{tabularx}
\end{table}

\subsection{Carrier atlas and C1--C7 summary}

The carrier atlas records the surfaces that matter for the final interpretation. The
coarse public carrier, \nolinkurl{intrinsic_naked_interface_coarse}, is kept as a
public signature shadow with \texttt{signature\_only} status. The lawful same-branch
witness, \nolinkurl{intrinsic_naked_interface_witness_rewrite_effective_operator}, is
the earned lawful survivor. The composite basis
\nolinkurl{witness_site_public_shadow} is the operational two-level basis: usable,
fragile, not a strongest new carrier in its own right. The scoped protocol-indexed
carrier, \nolinkurl{intrinsic_naked_interface_witness_autonomous_protocol_lift}, is
the strongest survived positive carrier surface, at flagship-only / protocol-indexed
scope.

The atlas also records the surfaces that block overclaim. The nearby-positive panel
is a real scope-limitation row; the support-only erasure-recovery anchor is positive
only as support-only transfer; the global/full-core overclaim row remains collapsed
under honesty; and the adjacent-surface search remains at no exactification.

At the C1--C7 level the atlas condenses the same picture. The coarse public row has
partial readout and transport, fails C3, and fails C7. The lawful same-branch witness
passes C1--C6 and is only partial at C7, which is why it is a real same-branch
survivor but not a cross-seed/full-core carrier. The operational composite is usable
but fragile. The scoped protocol-indexed carrier passes section honesty, is partial
on readout, transport, protocol honesty, and stress/expansion, and fails both C3 and
C7 for strongest-new-carrier purposes. The nearby-positive and global/full-core rows
stay negative.

Table~\ref{tab:carrier-atlas-c1c7} gives the compact form.

\begin{table}[t]
\centering
\scriptsize
\caption{Carrier atlas and compact C1--C7 matrix. The atlas preserves the distinction
between the lawful same-branch witness, the public signature shadow, and the scoped
protocol-indexed survivor.}
\label{tab:carrier-atlas-c1c7}
\begin{tabularx}{\linewidth}{@{}>{\raggedright\arraybackslash}p{0.20\linewidth}*{7}{>{\raggedright\arraybackslash}p{0.06\linewidth}}>{\raggedright\arraybackslash}X@{}}
\toprule
Atlas row & C1 & C2 & C3 & C4 & C5 & C6 & C7 & Final interpretation \\
\midrule
Public shadow (\nolinkurl{intrinsic_naked_interface_coarse}) & part & part & fail & part & part & part & fail & public-facing signature shadow only \\
Lawful witness (\nolinkurl{intrinsic_naked_interface_witness_rewrite_effective_operator}) & pass & pass & pass & pass & pass & pass & part & real same-branch lawful survivor \\
Operational basis (\nolinkurl{witness_site_public_shadow}) & part & part & fail & part & part & part & part & usable two-level basis, but fragile and not strongest \\
Scoped carrier (\nolinkurl{intrinsic_naked_interface_witness_autonomous_protocol_lift}) & part & part & fail & pass & part & part & fail & strongest survived positive carrier surface, but indexed and flagship-only \\
Nearby-positive panel & fail & fail & n/a & fail & fail & fail & fail & real scope limitations block broader transfer \\
Support-only anchor & fail & fail & n/a & pass & pass & part & fail & support-only positive row, not carrier exactification \\
Global/full-core overclaim & fail & fail & fail & fail & fail & fail & fail & not licensed \\
Adjacent-surface search & fail & fail & n/a & n/a & n/a & fail & fail & no upgraded new surface earned \\
\bottomrule
\end{tabularx}
\end{table}

\subsection{The C3 non-reducibility audit}

C3 is the criterion that most directly explains why the verdict is scoped. The gate
compares the richest surviving candidate against the lower reduction targets of
Section~\ref{sec:candidate-classes-reductions-criteria}---\(Q\), \(M\), \(K\), \(R\),
and \(\Lambda^\ast\)---and against the already-landed lawful witness lane and the
witness/public-shadow basis.

For final-gate purposes the strongest surviving scoped carrier is
\nolinkurl{intrinsic_naked_interface_witness_autonomous_protocol_lift}. C3 records it
as admissible as a weaker positive surface but not clearing the strongest new-carrier
audit: the object reduces back to the already-landed lawful witness lane. The dominant
reduction target is
\nolinkurl{intrinsic_naked_interface_witness_rewrite_effective_operator}. The scoped
carrier survives honestly; the gate is not licensed to promote it into a strongest
new-carrier verdict independent of the simpler lawful witness lane it enriches.

The scoped carrier is the flagship-only protocol-indexed surface on which the
quadratic branch carries three claim-bearing quadratics and the proto-event pilot
finds two proto-event/test families. C3 does not erase those positives; it classifies
them as evidence that a weaker, scoped carrier surface survives, not that a new
nonindexed/full-core carrier has overtaken the same-branch witness.

The same C3 matrix records the primary lawful witness as
\texttt{resolved\_nonreducible}: no source-first reduction witness collapses
\nolinkurl{intrinsic_naked_interface_witness_rewrite_effective_operator} into a
simpler descriptor or process baseline. C3 therefore preserves the lawful same-branch
witness as a real positive survivor while blocking the strongest carrier claim from
migrating to the richer-looking but still reducible protocol-indexed surface.

The operational witness/public-shadow basis is useful but reducible to the lawful
witness lane; the coarse public carrier is reducible to the weaker public component
of that basis. The final gate therefore distinguishes three kinds of positive or
quasi-positive object: a lawful witness that lands, an operational basis that is
useful but not strongest, and a scoped carrier surface that survives but reduces for
strongest-claim purposes. This is the content of
\texttt{protocol\_indexed\_carrier\_only}.

\subsection{Claim permissions}

With the atlas and C3 audit fixed, the claim-permission map is straightforward. The
paper may claim a lawful same-branch witness; a flagship-only, protocol-indexed
carrier surface; and scoped quadratic carrying and scoped proto-event/test families on
that same survived carrier. It may not claim law on the public carrier, cross-seed or
full-core carrier exactification, a general quadratic or covariance bridge, a general
event geometry, pathwise filtering structure, or loop-action success.

Table~\ref{tab:claim-permissions} gives the compact permission map.

\begin{table}[t]
\centering
\footnotesize
\caption{Claim permissions after the final gate. Positive permissions remain scoped;
stronger public, cross-seed, and general-geometry claims are not licensed.}
\label{tab:claim-permissions}
\begin{tabularx}{\linewidth}{@{}>{\raggedright\arraybackslash}p{0.34\linewidth}>{\raggedright\arraybackslash}p{0.14\linewidth}>{\raggedright\arraybackslash}X@{}}
\toprule
Claim row & Permission & Scope / note \\
\midrule
Lawful same-branch witness & licensed & same-branch positive survivor \\
Scoped protocol-indexed carrier & licensed & flagship-only, protocol-indexed survivor \\
Law on the public carrier & not licensed & public row remains signature-only \\
Cross-seed carrier & not licensed & nearby-positive panel remains outside exact transfer \\
Full-core carrier & not licensed & full-core/global overclaim row remains collapsed \\
Scoped quadratic carrying & licensed & only on the survived flagship-only, protocol-indexed carrier \\
General quadratic or covariance bridge & not licensed & comparison-only bilinear row is not promotable \\
Scoped proto-event/test families & licensed & only on the same survived scoped carrier \\
General event geometry & not licensed & proto-event pilot remains selective and local to the scoped carrier \\
Pathwise filtering structure & not licensed & route filtering remains unavailable \\
Loop-action success & not licensed & loop-action lane remains not found \\
\bottomrule
\end{tabularx}
\end{table}

The secondary flags fit inside this permission map. The gate preserves the presence
of the lawful same-branch witness, the signature-only status of the coarse public
carrier, the non-faithful witness/public compression relation, cross-seed
partiality, the nearby positives as real scope limitations, the collapse of the
broader site-not-probe honesty claim, the absence of a new-surface upgrade, and the
scoped quadratic and proto-event rows. It also records several negative diagnostics
explicitly: no broader covariance bridge, route filtering unavailable, loop-action not
found, and the paired panel kept separate.

\subsection{Interpretation of the gate}

The strongest public, cross-seed, and full-core carrier claim is not licensed. The
lawful same-branch witness is licensed. The public signature shadow remains coarse
and non-faithfully related to that witness. A flagship-only, protocol-indexed carrier
surface is licensed as the strongest survived positive carrier row, with the
downstream quadratic and proto-event rows as secondary scoped structures on that
same survivor. The final gate is positive but restricted.

The scoped verdict is not a weakened version of a larger claim that almost landed. It
is the verdict for the artifact surface the repository earned. The lawful
witness/public-shadow split is central; the scoped carrier is real but indexed; the
nearby positives remain real scope limitations; the stress, honesty, and
adjacent-surface panels are part of the verdict; and the secondary quadratic and
proto-event rows confirm only that the scoped survivor is nontrivial, not that it
globalizes.

Section~\ref{sec:formal-and-lean-criterion-layer} turns to the formal layer: how the
Lean criterion kernel and witness-data bridge formalize the vocabulary, criteria, and
scoped witness data used to state the verdict.
 \section{Formal Criterion Layer and Lean Bridge}
\label{sec:formal-and-lean-criterion-layer}

By the end of Section~\ref{sec:final-gate-and-carrier-atlas}, the empirical carrier
gate is closed. The final verdict, the carrier atlas, the scope limitations, and the
claim permissions are artifact-backed results of the pipeline. The formal layer gives
a typed vocabulary for the carrier-exactification question, proves a kernel of
abstract gating lemmas, mirrors the final scoped verdict in Lean-readable witness
data, and separates those formal objects from the empirical rebuild surface.

The paper relies on two kinds of support. Lean formalizes the vocabulary and the
abstract implication structure of the gate; the canonical pipeline rebuilds or
validates the empirical artifact surface. Neither layer is a theorem that the other
is correct.

\subsection{Typed vocabulary for carrier exactification}

The formal layer is written in the Lean theorem prover \citep{deMouraEtAl2015}; the
namespace is \nolinkurl{Recombination.CarrierExactification}, and
\Cref{app:lean-formalization-summary} lists the file-level module layout. The
typed vocabulary used there instantiates the criterion-based operational approach
familiar from \citet{Hardy2001} and \citet{ChiribellaDArianoPerinotti2011}; no new
mathematics is claimed, and the Lean layer contributes typechecking of the abstract
implications rather than a proof of the empirical verdict.

We typeset the Lean definitions below in mathematical form while keeping the
Lean-side predicate reading explicit: the relations are \textup{Prop}-valued in
the code, and can be read set-theoretically in the usual way.

\begin{definition}[Designated surface]
A \emph{designated surface} is a tuple
$\Sigma = (\mathcal{S}, \mathcal{O}, \mathcal{G}, \Xi, \mathcal{P})$ of seed,
observable, generator, section, and protocol types.
\end{definition}

\begin{definition}[Carrier candidate]
\label{def:carrier-candidate}
A \emph{carrier candidate} on $\Sigma$ is a tuple
$\mathcal{C} = (X, V, \rho, \lambda, \tau)$ consisting of
\begin{itemize}[leftmargin=2em, itemsep=0pt, topsep=2pt]
\item a state space $X$,
\item a validity predicate $V \colon X \to \mathrm{Prop}$,
\item a seed-realization relation $\rho \colon \mathcal{S} \to X \to \mathrm{Prop}$,
\item a readout relation $\lambda \colon X \to \mathcal{O} \to \mathrm{Prop}$,
\item a transport relation $\tau \colon \mathcal{G} \to X \to X \to \mathrm{Prop}$.
\end{itemize}
\end{definition}

The seven criteria C1--C7 of
Section~\ref{sec:candidate-classes-reductions-criteria} are expressed as predicates
on $(\mathcal{C}, \Sigma)$, with auxiliary section, protocol, expansion, and
seed-link witness relations supplied where required. Two representative examples are
\begin{align*}
\operatorname{RS}(\mathcal{C}, \mathcal{S}_{\textup{des}}) &\;\iff\;
\forall s \in \mathcal{S}_{\textup{des}}:\ \exists x \in X,\ o \in \mathcal{O}:\
V(x) \wedge \rho(s, x) \wedge \lambda(x, o), \\
\operatorname{TC}(\mathcal{C}) &\;\iff\;
\forall g \in \mathcal{G},\ x, y \in X:\
V(x) \wedge \tau(g, x, y) \Rightarrow V(y),
\end{align*}
where $\operatorname{RS}$ is readout sufficiency (C1) on a designated seed scope
$\mathcal{S}_{\textup{des}} \subseteq \mathcal{S}$, and $\operatorname{TC}$ is
transport closure, one half of the law side of C2. The companion predicate
$\operatorname{RTC}(\mathcal{C})$ for readout/transport compatibility completes
C2 in the Lean kernel, while $\operatorname{TE}(\mathcal{C})$ is an additional
transport-exactness strengthening carried by the full evidence record. Section
honesty, protocol honesty, expansion honesty, and cross-seed support are
defined analogously from the corresponding Lean witnesses.

\begin{definition}[Reduction]
A \emph{reduction} from $\mathcal{C}$ to $\mathcal{D}$ on the same surface $\Sigma$
is a map $\phi \colon X_{\mathcal{C}} \to X_{\mathcal{D}}$ satisfying, for all
admissible arguments,
\[
V_{\mathcal{C}}(x) \Rightarrow V_{\mathcal{D}}(\phi(x)), \qquad
\rho_{\mathcal{C}}(s, x) \Rightarrow \rho_{\mathcal{D}}(s, \phi(x)),
\]
\[
\lambda_{\mathcal{C}}(x, o) \Rightarrow \lambda_{\mathcal{D}}(\phi(x), o),
\qquad
\tau_{\mathcal{C}}(g, x, y) \Rightarrow \tau_{\mathcal{D}}(g, \phi(x), \phi(y)).
\]
We write $\mathcal{C} \preceq \mathcal{D}$ when a reduction exists, and say
$\mathcal{C}$ is \emph{non-reducible} among a family
$\{\mathcal{D}_i\}_{i \in I}$ if $\mathcal{C} \not\preceq \mathcal{D}_i$ for every
$i$. This is the C3 comparator of
Section~\ref{sec:candidate-classes-reductions-criteria}.
\end{definition}

Scope, law, compression, indexing, quadratic, proto-event, and verdict statuses are
encoded as finite enumerations (for example, the law statuses $\{\text{lawful},
\text{signature-only}, \text{collapse}\}$, the compression statuses
$\{\text{faithful}, \text{non-faithful}, \text{n/a}\}$, and the verdict values
including $\text{protocol-indexed-carrier-only}$). Evidence bundles are structured
records that package the predicates above alongside a choice of scope, honesty and
expansion witnesses, and a baseline family for the C3 comparison; the three bundles
used below are \emph{scoped}, \emph{protocol-indexed}, and \emph{full
exactification} evidence. When the main text writes ``lawful witness,'' ``public
shadow,'' or ``protocol-indexed carrier,'' these phrases correspond to typed
distinctions on $\Sigma$; the definitions do not certify that a particular
JSON artifact has the right value.

\subsection{Criterion kernel}

The proof kernel file typechecks three groups of implications on top of the
vocabulary above. The local target
\nolinkurl{Recombination.CarrierExactification} compiles cleanly and carries no
active \texttt{sorry}. The use of a small theorem kernel inside a proof assistant
has clear precedent in large formalization efforts \citep{Gonthier2008}; our use is
narrower, since the implications below are
logical consequences of the definitions in
Section~\ref{sec:candidate-classes-reductions-criteria} rather than new theorems
about carrier theory.

\emph{Reduction interface.} The kernel establishes identity reduction,
composition, and reflexive/transitive reducibility, and records that reductions
preserve validity, readout, and transport. Non-reducibility $\mathcal{C}
\not\preceq \mathcal{D}_i$ for every $i$ is derived from the preservation laws.

\emph{Evidence extraction.} From a scoped-positive evidence bundle, the kernel
extracts scoped positivity and the criterion predicates covered by that bundle.
From a full-exactification bundle, it extracts readout sufficiency, transport
closure and exactness, section honesty, protocol honesty, expansion honesty,
cross-seed support, and structural non-reducibility against the bundle's
baseline family. The summary implication has the following shape.

\begin{proposition}[Criterion implication, abstract form]
\label{prop:criterion-implication}
Let $\Sigma$ be a designated surface and $\mathcal{C}$ a carrier candidate on
$\Sigma$, and let $E$ be a full carrier-exactification evidence bundle for
$\mathcal{C}$ with designated seed scope
$\mathcal{S}_{\textup{des}} \subseteq \mathcal{S}$, witness families for section,
protocol, and expansion honesty, a seed-link relation, and a baseline family
$\{\mathcal{D}_i\}$. Then $\mathcal{C}$ satisfies the paper's seven criteria on
$\Sigma$: C1 via $\operatorname{RS}(\mathcal{C}, \mathcal{S}_{\textup{des}})$; C2
via $\operatorname{TC}(\mathcal{C})$ together with
$\operatorname{RTC}(\mathcal{C})$; C3 via $\mathcal{C} \not\preceq \mathcal{D}_i$
for every $i$; C4 via $\operatorname{SH}(\mathcal{C})$; C5 via
$\operatorname{PH}(\mathcal{C})$; C6 via $\operatorname{EH}(\mathcal{C})$; and C7
via $\operatorname{CS}(\mathcal{C}, \mathcal{S}_{\textup{des}})$. In addition, the
full evidence bundle yields the stronger transport-exactness predicate
$\operatorname{TE}(\mathcal{C})$.
\end{proposition}

The Lean theorem package behind this summary consists of
\texttt{carrier\_exactified\_of\_evidence} together with the companion
extraction lemmas (\texttt{fullEvidence\_extracts\_readoutSufficient} etc.),
which provide the component projections used here.

\emph{Scope barriers.} Protocol-indexed evidence implies scoped positivity, and a
companion lemma blocks its re-interpretation as a nonindexed upgrade; barrier
lemmas record that full exactification does not follow if any of cross-seed
support, section honesty, or protocol honesty is universally blocked. A
faithful-compression implication formalized via a \texttt{CompressionContext}
object captures the witness/public-shadow logic of
Section~\ref{sec:same-branch-law-and-public-shadow} at criterion level.

\begin{remark}
Each of these implications is a logical consequence of the typed definitions and
evidence records summarized above and in
Section~\ref{sec:candidate-classes-reductions-criteria}; they are not new
empirical theorems. The kernel does not establish that the predicates are
instantiated by the ticket31 JSON surface; that bridge is artifact-backed and is
preserved as an explicit trust boundary in
\Cref{app:lean-formalization-summary}.
\end{remark}

\subsection{Witness-data bridge}

The witness-data bridge adds a third layer between abstract theorems and empirical
artifacts: a Lean-readable witness-data mirror of the final verdict surface. In
\texttt{WitnessData.lean}, the key
structure is \texttt{Ticket31WitnessData}, together with a small \texttt{ClaimPermissions}
record and a \texttt{CarrierLabel} wrapper for the named ticket31 objects. The resulting
instance, \texttt{ticket31FinalWitnessData}, does not attempt to reproduce the entire
artifact graph. Instead it records the stable final vocabulary the paper needs.

The instance mirrors the bounded outcome: final verdict
\texttt{protocolIndexedCarrierOnly}; coarse public carrier \texttt{signatureOnly};
same-branch witness \texttt{lawful}; surviving carrier at \texttt{flagshipOnly} scope
with \texttt{protocolIndexedOnly} indexing; compression status \texttt{nonFaithful};
quadratic and proto-event rows at scoped positive statuses; permission booleans false
for carrier law at the public surface, cross-seed carrier, full-core carrier, and any
broader bridge or general event-geometry permission. Sanity lemmas state these facts
propositionally.

The witness-data file is a typed mirror of the final vocabulary and permissions, not
a theorem that the JSON artifacts are correct. The trust boundary is recorded
explicitly: artifact correctness, empirical C1--C7 satisfaction, the empirical C3
audit, runner correctness, cargo/PICA correctness, and the correctness of the
extraction bridge remain external.

The witness-data bridge lets the manuscript say that the final verdict and claim
permissions have a typed Lean-facing representation. It does not let the manuscript
say that the final verdict has been proved inside Lean from the experiment outputs.

\subsection{Reproducibility surface}

The empirical side is handled by the canonical runner
\nolinkurl{scripts/ticket31/run_canonical_pipeline.py} and the machine-readable asset
index \nolinkurl{results/derived/ticket31_asset_index.json}. In the current repo
state, the runner records the pipeline as canonical-ready with known gaps and the
executed mode as \texttt{check-only}.
The empirical surface is organized around a single canonical runner, a serial phase
order, contract tests, and an asset index that records what exists, what is stable,
what is environment-dependent, and what was validated rather than rebuilt.

Serial discipline matters because several steps rewrite shared derived JSON. The
runner fixes an explicit phase order: seed corpus, live-surface exactification,
law/PICA bridge, strengthening lane, cross-seed and falsifier decision path, Lean
build and witness bridge, and final contract checks. The asset index records Lean
summary data, vendor/PICA summary data, determinism notes, and known gaps such as
validated-only legacy comparison lanes.

The empirical claims are supported by the repository's machine-readable artifacts,
canonical runner, regression checks, and final-contract layer. The contract layer
preserves that the scoped quadratic-carrying and proto-event results remain scoped to
the flagship-only, protocol-indexed survivor, and does not let the secondary positives
expand into a broader verdict.

Table~\ref{tab:formal-repro-layers} summarizes the division of labor between the Lean
criterion layer, the witness-data mirror, the canonical pipeline, and the remaining
external trust boundary.

\begin{table}[t]
\centering
\footnotesize
\caption{Formal and reproducibility layers. Lean formalizes the carrier-exactification
vocabulary and abstract gates; the canonical pipeline and asset index provide the empirical
rebuild surface.}
\label{tab:formal-repro-layers}
\begin{tabularx}{\linewidth}{@{}>{\raggedright\arraybackslash}p{0.16\linewidth}>{\raggedright\arraybackslash}p{0.28\linewidth}>{\raggedright\arraybackslash}p{0.24\linewidth}>{\raggedright\arraybackslash}X@{}}
\toprule
Layer & What is formalized or validated & What remains external & Role in the paper \\
\midrule
Lean definitions (L01) & designated surfaces, carrier candidates, reducibility, honesty predicates, exactification evidence, scoped and protocol-indexed evidence & empirical instantiation of those predicates & stable formal vocabulary for the manuscript's object distinctions \\
Criterion kernel (L02) & reduction lemmas, evidence extraction, abstract criterion theorem, blocking lemmas, abstract non-faithful compression implication & proof that ticket31 artifacts satisfy the premises & theorem-level statement of the gate's abstract logic \\
Witness-data bridge (L03) & typed mirror of the final verdict, permissions, and scoped-positive flags & JSON correctness, extraction correctness, empirical C1--C7 and C3 audits & Lean-readable summary of the final bounded verdict \\
Canonical pipeline (T20) & serial runner, asset index, contract checks, determinism notes, Lean validation hooks & environment-dependent runtime details and known-gap steps recorded in the asset index & reproducible empirical rebuild and validation surface \\
Remaining trust boundary & explicit proof-gap ledger and external assumptions & final soundness from artifacts to Lean, cargo/PICA execution, phase runners & keeps the paper from overstating what has been formally derived \\
\bottomrule
\end{tabularx}
\end{table}

\subsection{Trust boundaries and proof gaps}

\texttt{ProofGaps.lean} records the unfinished bridges. The main named gaps are the
empirical C1--C7 bridge, final-gate soundness from the JSON surface, the empirical C3
non-reducibility bridge, the JSON extraction trust boundary, and the cargo/PICA
runner trust boundary. Other gaps cover the same territory from nearby angles:
protocol-indexed non-upgrade semantics are only partially internalized, scoped
quadratic and proto-event structure is not yet attached to a richer higher-order
formal interface, and the witness-data layer remains a mirrored fact layer rather
than a fully checked empirical certificate.

The paper relies on the correctness of the JSON artifacts, the Python bridge and
extraction scripts, cargo/PICA execution, the phase runners and other pipeline steps,
the empirical C1--C7 audit, and the empirical C3 non-reducibility comparison. None of
those are proved inside Lean; they are made reproducible and inspectable through the
canonical runner, the asset index, and the artifact summaries. The accompanying
trust-boundary reasoning follows standard accounts of what formal verification can
and cannot guarantee in mathematical practice \citep{AvigadHarrison2014}.

Sections~\ref{sec:live-surface-exactification}--\ref{sec:final-gate-and-carrier-atlas}
are empirical and artifact-backed. The formal layer supports their vocabulary and
implication structure, while the reproducibility layer supports their rebuild
discipline. Lean formalizes the vocabulary, the abstract gates, and a typed mirror of
the final scoped verdict; the canonical pipeline and asset index provide the
empirical rebuild surface; the empirical verdict itself remains artifact-backed and
reproducibility-checked rather than derived inside Lean.
 \section{Discussion}
\label{sec:discussion}

Paper~6 left a sharper question than whether the promoted package had interesting
internal structure: whether the surviving package exactified into a carrier, and at
what level. The preceding sections separated several objects that a looser account
would have collapsed, and the final gate fit them into one bounded verdict. This
section interprets the completed result.

\subsection{What exactified}

Two positive results sit at the center of the paper. A lawful same-branch witness
exactifies at the refined witness level: the decisive object is
\nolinkurl{intrinsic_naked_interface_witness_rewrite_effective_operator}, and
Section~\ref{sec:same-branch-law-and-public-shadow} showed that this witness, not the
coarse public surface, is where the same-branch law lands.

A carrier surface survives at the scoped protocol-indexed level: the surviving row is
\nolinkurl{intrinsic_naked_interface_witness_autonomous_protocol_lift}, with survival
class \nolinkurl{protocol_indexed_carrier_only} at \texttt{flagship\_only} scope. The
witness gives the law-bearing same-branch object; the protocol lift gives the
strongest surviving carrier surface retained by the final gate.

On the scoped carrier, three claim-bearing quadratic families remain carried and two
proto-event/test families remain present. These secondary positives show that the
scoped carrier is mathematically active; they do not replace the main result, and the
final gate records the whole outcome under the single bounded verdict
\texttt{protocol\_indexed\_carrier\_only}.

The promoted package does not remain only at the level of descriptor, process family,
chart, or support signal. It exactifies at the refined witness level and preserves a
distinct flagship-only protocol-indexed carrier surface above that witness.

\subsection{What did not exactify}

No lawful carrier lands on the public surface; no cross-seed carrier exactifies; no
full-core carrier exactifies; the nearby positives do not become exact transfers of
readout or law; the lifted flagship row does not promote to a nonindexed intrinsic
carrier; the adjacent-surface search does not produce an upgraded surface; and the
downstream secondary rows do not turn into a broader bridge or a general event
geometry.

The paper does not describe a strong carrier claim that almost lands and is saved by
cautious wording. It finds a lawful same-branch witness, a public signature shadow,
and a scoped protocol-indexed carrier surface. The strongest public, cross-seed, and
full-core claims are not licensed because the evidence surface is stratified.

In particular: the four nearby positives are not null cases or hidden positives; they
align at support level and remain real scope limitations. The support-only
erasure-recovery row is useful evidence as support-only transfer, not as exact
readout or lawful transport transfer. Controls remain controls.

The protocol-indexed survivor is not a disguise for nonindexed intrinsic promotion;
its indexing is part of the content of the result. The paper speaks of a carrier
surface that survives only with protocol coordinates attached, rather than a full
flagship carrier exactification with an incidental protocol annotation.

\subsection{The public signature shadow}

The public carrier is not a defective version of the lawful witness; it is a coarser
object with a different role. It records the public signature or readout surface of
the lawful witness without preserving the transport law that exactifies on the
refined same-branch witness. The refined witness plays, at carrier level, the role
that structural features play in ontic structural realism
\citep{FrenchLadyman2003}: the lawful content is located in the refined relational
structure rather than in the coarse public object.

\texttt{non\_faithful\_public\_compression} names the structural relation between the
two. The refined lawful witness is where the same-branch law lands; the coarse
public object remains \texttt{signature\_only}. The compression preserves a real
public role but does not preserve the full law.

The public shadow is the public-facing layer appropriate to the structure the run
actually earns. Public signature and readout behavior can remain stable even when the
full transport law lives only on a more refined witness. Two questions are therefore
separated: whether a lawful object exists, and whether the public surface itself
carries that law. The answers are yes and no, without collapsing into the
disappearance of public structure.

If the public shadow were discarded, the contrast that makes the witness result
meaningful would be lost. The result is neither ``law exists somewhere'' nor
``public readout exists somehow'' but the split itself: a lawful refined witness
together with a coarser public signature surface under non-faithful compression.

\subsection{Why the verdict is positive but scoped}

The final gate summarises the paper as ``strongest claim failed, survived positive
surface present.'' The two clauses answer different questions: whether the run earns
a new public, cross-seed, or full-core carrier verdict (no), and whether any honest
carrier surface remains after the reductions, honesty panels, stress panels, and
scope tests (yes).

\texttt{protocol\_indexed\_carrier\_only} is therefore not equivalent to
\texttt{no\_carrier\_exactified}. The lawful same-branch witness is a genuine
positive survivor, the autonomous protocol lift is a genuine scoped carrier surface,
and quadratic carrying and proto-event/test families survive on the same surface.

The verdict is also not \texttt{flagship\_only\_carrier\_exactified} in an unindexed
sense. The protocol-coordinate requirement remains binding at the gate, and C3
reduces the richest surviving carrier surface back to the already-landed lawful
witness lane for strongest-new-carrier purposes. The lifted surface cannot be
promoted into a new strongest carrier with indexing as an optional detail.

Scoped quadratic carrying and scoped proto-event/test families show that the
surviving carrier is nontrivial without broadening the carrier claim. Scope and
indexing are constitutive of the result, not caveats added afterward. A positive
result on a restricted scope is a familiar mode of progress in the physics of
emergence: coarse-grained laws can be genuine even when they neither replace nor
follow from the finer-grained ones \citep{Butterfield2011}.

\subsection{Limitations}

The limitations are structural. The surviving carrier remains flagship-only; protocol
indexing remains required; the public surface remains \texttt{signature\_only};
cross-seed reach remains partial, and the four nearby positives remain real scope
limitations rather than exact transfers; the broader site-not-probe claim fails the
honesty audit; route filtering remains unavailable; loop-action asymmetry is not
found; C3 blocks strongest-new-carrier language for the richest surviving scoped
carrier; no adjacent surface exactifies; and the scoped quadratic and proto-event
results do not constitute a general geometry.

These are the conditions under which the result is true, not corrections added
afterward. How public readout and test structure should be studied from such a
scoped carrier is a natural next question but is not taken up here.
 \section{Philosophical reading}
\label{sec:philosophical-reading}

The central asymmetry recorded in this paper admits a short philosophical
statement worth making explicit. Readout descends from the refined witness to
the public shadow; transport law does not. The two surfaces carry different
kinds of content, and the compression between them is structurally one-way.

In a Humean idiom, one may read that split as a distinction between how a
system appears at the public surface and how its law is carried at the refined
level. The public surface carries what the system looks like; the witness
carries how the system is governed to change. On that reading,
\emph{non-faithful public compression} is a typed way of making precise the
thought that the second cannot be recovered from the first by inspection alone:
the signature space does not determine the law space, because the compression is
many-to-one on state pairs that the law separates.

The witness/public-shadow split can also be read through the ontic--epistemic
contrast discussed by \citet{HarriganSpekkens2010} and sharpened, under further
assumptions, by \citet{PuseyBarrettRudolph2012}. The public shadow plays the
epistemic role in that comparison: several refined witness states can share the
same public image. The refined witness plays the ontic-role analogue: it
carries the content the transport law actually uses. The point is not that
Paper~7 proves an ontological thesis, but that its enrichment ceiling blocks a
reconstruction of the lawful carrier from public readout alone. What is lost
under compression is therefore not merely fine-grained state detail but the
law-bearing structure itself.

The atlas likewise admits a structural-realist gloss in the sense of
\citet{FrenchLadyman2003}. What one would naively call \emph{the object} of the
theory---the public carrier, the thing with a stable readout signature---is the
coarser, non-law-bearing projection. The lawful content instead lives in the
refined relational structure on the witness side. The lesson is not that the
public shadow is unreal, but that the law-bearing structure is explanatorily
prior to the coarse public object if one asks where the carrier-level law
actually resides.

Finally, the substrate exhibits a mode of partial descent that does not fit
cleanly either a reductionist picture or a strong-emergence picture. Readout
descent succeeds: downward determination of signatures holds. Transport-law
descent fails: downward determination of dynamics does not. This is a layered
result within a single atlas, with readout and law ending up on different
surfaces. \emph{Non-faithful public compression} is the paper's name for that
partial descent.

None of these readings is claimed as a theorem; they are philosophical
interpretations of a structural result that is itself artifact-backed and
criterion-formalized. The point of making them explicit is that the atlas
contributed by this paper admits a reading that connects to a longer-standing
set of questions---about phenomena and necessity, about ontic structure
under epistemic projection, about what counts as \emph{the object} in a
layered theory---rather than remaining only a scope-restricted technical
verdict.
 \section{Conclusion}
\label{sec:conclusion}

Paper~6 left the carrier status of the promoted survivor open. The present paper
contributes a carrier-level atlas on the same fixed-support Route~A flagship corpus.

The atlas consists of a \emph{lawful same-branch witness},
\nolinkurl{intrinsic_naked_interface_witness_rewrite_effective_operator}, on which
the refined same-branch transport and readout law exactifies; a \emph{coarse public
signature shadow}, \nolinkurl{intrinsic_naked_interface_coarse}, which retains
public signature role without carrying the refined lawful structure; and a
\emph{flagship-only protocol-indexed carrier surface},
\nolinkurl{intrinsic_naked_interface_witness_autonomous_protocol_lift}, built from
an autonomous lift of the witness. The relation between the witness and the public
shadow is recorded as \nolinkurl{non_faithful_public_compression}: a projection
preserving public signature role but not the refined lawful structure.

These four items are the objects we intend the series to carry forward. The
seven-criterion exactification audit adapted from operational-theory methodology,
together with denominator and expansion stress panels, a section-and-protocol
honesty rerun on the lifted flagship patch, and a controlled adjacent-surface
search, records that the witness passes all seven criteria at same-branch scope,
the public shadow passes only signature-level surface, and the protocol-indexed
carrier passes at flagship-only indexed scope. The final gate records the joint
classification as \texttt{protocol\_indexed\_carrier\_only}, with scope and
indexing preserved as part of the substrate rather than as caveats.

Three claim-bearing quadratic families and two proto-event/test families survive on
the flagship-only protocol-indexed carrier as scoped secondary structure. Carrier
law on the public surface, cross-seed carrier exactification, full-core carrier
exactification, and a broader event geometry are not earned; cross-seed
canonicalization remains partial and four nearby positives remain real scope
limitations. These non-claims are preserved explicitly in the claim-permission map
and mirrored in the Lean witness-data record.

The carrier-level layer the series carries forward after this paper is therefore
the refined lawful witness, the coarse public signature shadow, the protocol-indexed
carrier surface above them, and the non-faithful compression between witness and
shadow. That atlas, rather than the Paper~6 promoted-package layer alone, is the
relevant starting substrate after the present result.
 
\appendix
\section{Formal Definitions and Criteria}
\label{app:formal-definitions-and-theorem-statements}

This appendix expands the conceptual setup of
\Cref{sec:candidate-classes-reductions-criteria} without repeating the main text's
carrier atlas or final-gate tables. Its role is to fix the paper's formal vocabulary at a
prose-and-math level. The appendix does not introduce any stronger empirical claim than
the main body. In particular, the verdict vocabulary remains the ticket31 vocabulary:
lawful same-branch witness, coarse public signature shadow, scoped protocol-indexed
carrier survival, real scope limitations beyond the flagship anchors, and the final gate
\texttt{protocol\_indexed\_carrier\_only}.

\subsection{Designated surfaces and candidates}

The paper's abstract carrier question is always asked relative to a \emph{designated
surface}. A designated surface fixes the seed space, the observable space, the generator
or update family, and the section and protocol coordinates with respect to which honesty
questions can later be asked. A \emph{carrier candidate} on such a surface consists of a
state space \(X\), a validity predicate on \(X\), a realization relation assigning seeds
to valid states, a readout relation from states to observables, and a transport relation
describing admissible updates. The manuscript uses this common template to keep several
otherwise easy-to-confuse objects separate.

At the lowest level sit descriptor and process baselines. These include the current or
public quotient \(Q\), the predictive quotient \(M\) with factor map \(\pi : M \to Q\),
the branchwise quotient \(K\), the recombination quotient \(R\) with factor map
\(\eta : R \to K\), and the minimal behavioral carrier baseline \(\Lambda^\ast\).
Descriptor or process objects summarize lower behavior or runtime organization but do not
by themselves constitute a newly exactified carrier law. A chart-only intrinsic object may
still supply useful local coordinates, yet remain dependent on a local chart choice.
Section-mediated and section-fiber candidates introduce more structure, but still need to
clear the exactification criteria before they can be read as carrier-level survivors.

The main manuscript then distinguishes three rows that remain separate through the whole
paper. The first is the lawful same-branch witness, which is the first row that carries
law at the refined witness surface. The second is the coarse public carrier, which is kept
only as a public signature shadow. The third is the flagship-only protocol-indexed
survivor, which is a real carrier surface but only at the lifted indexed scope. The paper
therefore does not use ``carrier'' as a single undifferentiated noun. It uses it as a
typed family of claims whose scope, law status, and reducibility status all matter.

\subsection{Reduction and non-reducibility}

Reduction is the criterion that stops the manuscript from treating every richer-looking
object as a stronger result. A candidate \(C\) is \emph{reducible} to a baseline or
simpler positive lane \(B\) when there is an explicit map from the states of \(C\) to the
states of \(B\) preserving validity, seed realization, readout, and transport. In the
ticket31 setup, the first reduction targets are the lower baselines
\(Q, M, K, R,\Lambda^\ast\). The second tier consists of simpler positive rows already
earned in the run, especially the same-branch lawful witness and the two-level
witness/public-shadow basis.

This makes the C3 non-reducibility test concrete. A candidate fails C3 for a strongest
carrier claim not because it is ``insufficiently complicated,'' but because it adds no
claim-bearing state/readout/transport role beyond a simpler baseline or a simpler already
landed positive lane. Conversely, a candidate clears C3 only when no such reduction is
available among the relevant comparison rows. This is precisely why the final gate can
retain a scoped surviving carrier surface while refusing a strongest new-carrier verdict.
The scoped protocol-indexed lift is scientifically real, but for strongest-claim purposes
it still reduces back to the already-landed lawful witness lane.

The same distinction also clarifies the public shadow. The coarse public row is not merely
a weaker law row; it is a deliberately coarser surface that loses some of the structure
carried by the refined witness. That loss is recorded in the paper as
\texttt{non\_faithful\_public\_compression}. So reduction and compression do different
jobs. Reduction compares a candidate to simpler baselines or already-landed positives.
Compression compares the refined witness to the public readout shadow and explains why
public signature may remain meaningful without inheriting full law.

\subsection{Exactification criteria}

The main text introduces the seven exactification criteria C1--C7. Here they are stated in
slightly more explicit form.

\begin{enumerate}[label=\textbf{C\arabic*.}, leftmargin=*]
\item \textbf{Readout sufficiency.} The candidate must realize the designated seeds by
valid states with enough readout support to make the surface scientifically legible on the
intended panel. A support-only or signature-only row can remain useful without fully
passing this criterion.

\item \textbf{Transport or update law.} The candidate must support a stable update
relation on its state space. At strongest scope this means closure or exactness rather
than signature-level bookkeeping. Signature-only transport can support a live candidate or
public shadow without yielding full law.

\item \textbf{Non-reducibility.} The candidate must not collapse to the lower baselines or
to a simpler positive lane for the claim being made. This is the criterion that blocks the
richest surviving scoped carrier from being promoted to a strongest new-carrier result in
ticket31.

\item \textbf{Section honesty.} The candidate must survive admissible section changes on
the scope being claimed. A row can remain same-branch or flagship-robust while still
failing a broader section claim.

\item \textbf{Protocol honesty.} The candidate must survive the protocol-internalized
audit appropriate to the claim. A protocol-indexed survivor can pass this criterion at
indexed scope while still failing any nonindexed promotion claim.

\item \textbf{Stress and expansion honesty.} The candidate must remain scientifically
stable under denominator stress and repo-native surface expansion. Ticket31 uses this
criterion to keep coarse-public weakness, full-core overclaim, and adjacent-surface
overreach from being silently ignored.

\item \textbf{Cross-seed support.} The candidate must survive canonicalization beyond the
flagship anchors before a cross-seed or full-core claim is licensed. Partial support
alignment is meaningful, but it is not enough to promote the result into cross-seed or
full-core exactification.
\end{enumerate}

The criteria are therefore asymmetric by design. Failure of C7 blocks a cross-seed or
full-core claim without erasing a same-branch or flagship-only positive. Failure of public
law does not erase the lawful witness. Passing section and protocol honesty at scoped
protocol-indexed level does not imply nonindexed promotion. The paper's final verdict is a
direct consequence of this asymmetry.

\subsection{Gate vocabulary}

The gate vocabulary used in \Cref{sec:final-gate-and-carrier-atlas} is best understood as
a controlled verdict ladder rather than as a single scalar score. At the law layer the
important statuses are \emph{lawful}, \emph{signature-only}, and \emph{collapse}. At the
compression layer the important statuses are \emph{faithful}, \emph{non-faithful}, and
\emph{not applicable}; the present paper uses only
\texttt{non\_faithful\_public\_compression} for the witness-to-public relation. At the
scope layer the main manuscript distinguishes \texttt{flagship\_only},
\texttt{cross\_seed}, \texttt{full\_core}, \texttt{support\_only}, and
\texttt{new\_surface\_candidate} rows.

The verdict layer is intentionally richer than ``carrier'' versus ``no carrier.'' Full
exactification, chart-only survival, descriptor/process-only survival, blocked lanes, and
protocol-indexed-only survival all belong to the same vocabulary. Ticket31 closes with the
verdict \texttt{protocol\_indexed\_carrier\_only}. That verdict should be read together
with the lawful witness row, the signature-only public shadow, the nearby-positive scope
limitations, and the scoped secondary flags. The gate vocabulary therefore preserves two
different kinds of positive statement at once: exactification at the refined witness level
and scoped survival at the lifted protocol-indexed level.

This is also why the appendix keeps the public shadow explicit. The public row is not the
place where law exactifies, but it is still part of the paper's object vocabulary because
it records what remains visible on the public/readout surface. The appendices therefore
support the main text's two-level interpretation rather than replacing it with a simpler
story.
 \section{Computational Methods and Artifact Discipline}
\label{app:computational-methods-and-artifact-discipline}

This appendix expands the evidence-surface discussion of
\Cref{sec:corpus-artifacts-pipeline}. The main text keeps the computational story compact so
that the paper remains readable as a results paper rather than as a build log. Here the
focus shifts from narrative flow to artifact provenance. The appendix still remains
selective: it names the exact builders and summary artifacts that anchor the manuscript,
but it does not reproduce raw JSON or an exhaustive asset list. The exhaustive machine map
belongs instead to the canonical asset index in \Cref{app:reproducibility-and-source-map}
and the extended tables in \Cref{app:extended-tables}.

\subsection{Source recovery and seed freeze}

The computational starting point is the frozen seed corpus
\nolinkurl{results/derived/ticket31_seed_corpus.json}. That artifact records the fixed
eleven-row comparison surface: two flagship anchors, four nearby positives, four controls,
and one support-only row. In ticket31 this corpus is paired with the surface registry
\nolinkurl{results/derived/ticket31_surface_registry.json}, which records which operative
surfaces are builder-defined, which remain source-defined, and which legacy ticket30
artifacts are absent from the live snapshot.

This source-recovery step matters because ticket31 does not assume that every relevant
ticket30 payload is already materialized in the current worktree. Instead, the paper
rebuilds its auditable surface from repo-native sources: benchmark-backed config
references, relative-cycle parameterizations, paired artifact-control logic, and the code
or tests that define the supported readout, operator, section, protocol, and extension
surfaces. The builder entry points
\nolinkurl{scripts/ticket31/build_seed_corpus.py} and
\nolinkurl{scripts/ticket31/build_surface_registry.py} are therefore not mere convenience
scripts. They are part of the paper's evidence discipline: they freeze the corpus and
record what can be audited directly from the current repository.

\subsection{Baseline adapters}

Once the seed surface is frozen, ticket31 resolves the lower baselines through
\nolinkurl{results/derived/ticket31_reduction_baseline_smoke.json}. That artifact records
three resolution modes. Six rows are resolved by direct benchmark configuration, including
the flagship anchors, the ordinary controls, and the support-only row. Four nearby
positives are resolved as synthetic relative-cycle cases, because they are defined by a
parametric shortlist rather than by one pre-materialized benchmark bundle. The artifact
control is resolved as a logical paired control and remains \emph{paired-only} rather than
being flattened into an ordinary seed row.

For every resolved row the baseline layer records the availability of \(Q\), \(M\),
\(\pi\), \(K\), \(R\), \(\eta\), and \(\Lambda^\ast\). Ten rows are fully usable as
ordinary executable bundles; one row is explicitly usable only as a paired logical panel.
This is the concrete computational realization of the reduction targets introduced in
\Cref{sec:candidate-classes-reductions-criteria}. Baseline comparison in the paper is
therefore not a retrospective slogan. It is a machine-written surface over the same frozen
seed set as the later carrier candidates.

\subsection{Candidate and reduction artifacts}

The candidate layer is anchored by two summary artifacts:
\nolinkurl{results/derived/ticket31_candidate_census.json} and
\nolinkurl{results/derived/ticket31_reduction_map_summary.json}. The candidate census fixes
which object classes exist in the run at all: descriptor partitions, a process-object
family, a chart-only intrinsic side branch, section-mediated and section-fiber families, a
protocol-indexed object lane, the intrinsic carrier candidate, and the prospective
new-surface family. The reduction map then classifies these candidates as
descriptor/process baselines, chart-only rows, live carrier candidates, or unresolved
lanes relative to the lower baselines and the available support surfaces.

The live-surface stage adds three more artifacts:
\nolinkurl{results/derived/ticket31_intrinsic_live_surface_object.json},
\nolinkurl{results/derived/ticket31_live_surface_exactification_summary.json}, and
\nolinkurl{results/derived/ticket31_live_surface_branch_comparison.json}. Together these
artifacts build and audit the seed-tagged intrinsic object, record its direct-core and
paired-panel support, and compare it against the sectioned, section-fiber, and early
protocol-indexed lanes. That is the point at which the intrinsic branch becomes a real live
carrier candidate rather than a merely conceptual candidate class.

\subsection{PICA-backed support artifacts}

The later law and strengthening sections rely on a second layer of artifact support that is
more platform-specific but still machine-readable. The most important summary artifacts are
\nolinkurl{results/derived/ticket31_transport_law_summary_pica.json},
\nolinkurl{results/derived/ticket31_readout_law_summary_pica.json},
\nolinkurl{results/derived/ticket31_public_carrier_cutover_decision.json}, and
\nolinkurl{results/derived/ticket31_intrinsic_promotion_gate.json}. These artifacts do not
stand alone. They are tied to the PICA bridge and campaign slice layer, which includes
materials such as
\nolinkurl{results/derived/ticket31_pica_bridge_source_index.json},
\nolinkurl{results/derived/ticket31_pica_campaign_export.json}, and
\nolinkurl{results/derived/ticket31_pica_packaging_selection_ledger.json}.

The manuscript uses these objects in a limited way. They supply the lawful same-branch
witness ladder, the public-shadow cutover decision, the autonomous protocol lift, the
joint attack matrix, and the intrinsic promotion gate. They do not convert the paper into
a platform log. The paper reads them only through scoped summary claims: witness law,
signature-only public shadow, flagship-only protocol-indexed survival, and the explicit
limits carried into T12--T19.

\subsection{Stress, honesty, and scoped secondary structures}

The downstream falsifier and secondary-structure lane is likewise artifact-backed rather
than narrative-only. Cross-seed and scope-contract detail are recorded in
\nolinkurl{results/derived/ticket31_cross_seed_canonicalization_summary.json} and
\nolinkurl{results/derived/ticket31_nonflagship_scope_contract_summary.json}. Denominator
stress, surface expansion, and the refreshed section/protocol honesty gate are recorded in
\nolinkurl{results/derived/ticket31_denominator_stress_summary.json},
\nolinkurl{results/derived/ticket31_surface_expansion_summary.json}, and
\nolinkurl{results/derived/ticket31_section_protocol_gate.json}. The adjacent-surface search
is recorded in
\nolinkurl{results/derived/ticket31_new_surface_exactification_summary.json}. The scoped
quadratic and proto-event rows are recorded in
\nolinkurl{results/derived/ticket31_quadratic_carrying_summary.json} and
\nolinkurl{results/derived/ticket31_proto_event_pilot_summary.json}.

The computational discipline of these later artifacts matches the discipline of the earlier
ones. Nearby positives remain explicit as nearby positives rather than silently joining the
anchor panel. Controls remain controls and the paired artifact row remains paired. The
support-only anchor remains support-only. Scoped secondary structure is attached only to the
same flagship-only, protocol-indexed survivor recorded by the promotion gate, and the
adjacent-surface search remains explicit when it earns no exactification.

The result is a paper whose technical support is inspectable without becoming unreadable.
Each section in the main text relies on a small set of canonical summary artifacts. The
appendices name those artifacts and explain how they were produced. The repository remains
the evidence surface, but the appendices keep that surface selective enough to support the
main argument rather than displacing it.
 \section{Reproducibility and Source Map}
\label{app:reproducibility-and-source-map}

This appendix expands the reproducibility summary from
\Cref{sec:formal-and-lean-criterion-layer}. Its subject is the empirical rebuild surface,
not the carrier verdict itself. The paper's claims are anchored to machine-readable
artifacts, a serial canonical runner, and contract checks. They are not anchored to the
idea that every file in \nolinkurl{results/derived/} must be read manually. The goal of
this appendix is therefore to document the canonical rebuild path and the asset map in a
way that lets the reader trace the evidence surface without creating a second manuscript
inside the appendices.

\subsection{Canonical runner}

Ticket31's canonical entry point is
\nolinkurl{scripts/ticket31/run_canonical_pipeline.py}. The runner exposes both
\texttt{--check-only} and \texttt{--full} modes. T20 used \texttt{check-only}, which
validated the canonical artifact set, Lean target, and final contract without executing a
fresh full rebuild in that turn. The existence of a \texttt{--full} path matters for the
paper because it establishes that the serial rebuild plan is not merely descriptive. At
the same time, the current ticket31 reproducibility note is explicit that the T20 record is
``canonical ready with known gaps'' rather than a claim of perfect or gap-free rebuild.

The canonical phase order is fixed in the runner itself:
\begin{enumerate}[leftmargin=*]
\item \texttt{phase0\_preflight}
\item \texttt{phase1\_seed\_surface\_baselines}
\item \texttt{phase2\_live\_surface\_exactification}
\item \texttt{phase3\_generic\_law\_and\_pica\_bridge}
\item \texttt{phase4\_phase\_runners}
\item \texttt{phase5\_t11m\_strengthening\_lane}
\item \texttt{phase6\_post\_t11m\_decision\_lane}
\item \texttt{phase7\_lean\_bridge\_validation}
\item \texttt{phase8\_asset\_index\_and\_contracts}
\end{enumerate}
This ordering matters because several steps rewrite shared derived JSON outputs. The paper
therefore treats phase order as part of the evidence surface, not as a performance detail.

\subsection{Serial execution}

The runner is deliberately serial. Shared derived JSON writers make concurrent rebuilds
scientifically ambiguous even when they are computationally possible. The reproducibility
surface therefore uses serial execution as a correctness discipline rather than as a
convenience. The paired artifact row, the PICA-backed support artifacts, and the final
contract reports all rely on having one declared phase order and one declared action per
step.

This is also why the pipeline distinguishes \emph{executed}, \emph{validated}, and
\emph{validated-only} rows. The paper does not hide which surfaces were rebuilt in the
current turn and which were accepted as already-present canonical artifacts under
\texttt{check-only} mode. That distinction is part of the provenance of the empirical
surface, not a caveat to be moved out of sight.

\subsection{Asset index}

The central source map is
\nolinkurl{results/derived/ticket31_asset_index.json}. In the current repo state it records
\(190\) artifact entries together with the canonical phase order, counts by artifact kind,
required-status counts, determinism notes, Lean summary data, vendor/PICA summary data,
final-contract status, and explicit known gaps. For the paper, the most important point is
that the asset index is organized around the paper's actual claim surface rather than
around a generic filesystem crawl.

Several exact artifact paths are central enough to list explicitly here:
\begin{itemize}[leftmargin=*]
\item \nolinkurl{results/derived/ticket31_gate_decision.json}
\item \nolinkurl{results/derived/ticket31_carrier_type_atlas.json}
\item \nolinkurl{results/derived/ticket31_asset_index.json}
\item \nolinkurl{results/derived/ticket31_quadratic_carrying_summary.json}
\item \nolinkurl{results/derived/ticket31_proto_event_pilot_summary.json}
\item \nolinkurl{results/derived/ticket31_lean_witness_bridge.json}
\end{itemize}

These paths are not listed because they are the only relevant artifacts. They are listed
because together they show the shape of the paper's evidence surface: final gate,
carrier atlas, asset inventory, scoped secondary-structure summaries, and the Lean-facing
bridge artifact.

\subsection{Determinism and known gaps}

The T20 status is \nolinkurl{canonical_ready_with_known_gaps}. The determinism report is
\nolinkurl{results/derived/ticket31_pipeline_determinism_report.json}. It records
\(181\) stable artifacts, \(3\) environment-dependent artifacts, and \(6\) missing optional
artifacts. The known-gap ledger also records two intentionally absent standalone scripts:
\nolinkurl{scripts/ticket31/build_comparison_exactification.py} and
\nolinkurl{scripts/ticket31/build_transport_readout_law_rerun.py}. In both cases the
runner uses a module entry point or validates the already-present canonical artifact set
instead of pretending that a dedicated standalone builder still exists.

The other named known gap is that T20 did not execute a full rebuild in the current turn.
That point should not be hidden because it tells the reader exactly what the current
reproducibility statement means: the canonical artifact set and final contract were checked
successfully, but the historical ticket31 surface was not regenerated from scratch in that
specific run. The distinction between executed and validated-only surfaces follows the
computational-reproducibility recommendations of \citet{StoddenEtAl2016}.

\subsection{Final contract checks}

The canonical runner ends by writing
\nolinkurl{results/derived/ticket31_artifact_contract_report.json} and
\nolinkurl{results/derived/ticket31_final_contract_report.json}. Both have
\texttt{passed} status in the current snapshot. The artifact-contract layer checks the
branch decisions and scope guards that matter for the paper's integrity: partial
cross-seed status, real scope limitations on the nearby positives, null-by-design controls,
support-only handling, fragile denominator and expansion panels, honesty collapse of the
global/full claim, absence of adjacent-surface exactification, scoped T17/T18 status, and
the final T19 gate.

The final-contract layer adds a narrower set of manuscript-facing guards. It preserves the
final verdict \texttt{protocol\_indexed\_carrier\_only}; keeps the coarse public row at
\texttt{signature\_only}; keeps the lawful witness separate and \texttt{lawful}; keeps the
scoped carrier at flagship-only / protocol-indexed scope; keeps the paired panel explicit;
keeps route filtering unavailable; and keeps loop-action asymmetry
\texttt{not\_found}. This contract layer is the reason the appendices can remain
artifact-rich without silently widening the paper's claim surface.
 \section{Lean Formalization Summary}
\label{app:lean-formalization-summary}

This appendix expands the formal summary of
\Cref{sec:formal-and-lean-criterion-layer}. Its purpose is to describe the local Lean
layer precisely enough that the reader can see what is formalized, what is mirrored into
Lean-readable data, and what remains external. The appendix does \emph{not} turn the paper
into a formal-certification claim. The empirical verdict remains artifact-backed and
reproducibility-checked. The Lean layer supports vocabulary, abstract gate structure, and a
typed mirror of the final bounded verdict.

\subsection{Lean module layout}

The relevant local namespace is \texttt{Recombination.CarrierExactification}. The root
module is \nolinkurl{lean/Recombination/CarrierExactification.lean}. The internal layout is
split across:
\begin{itemize}[leftmargin=*]
\item \nolinkurl{lean/Recombination/CarrierExactification/Definitions.lean}
\item \nolinkurl{lean/Recombination/CarrierExactification/Statements.lean}
\item \nolinkurl{lean/Recombination/CarrierExactification/ProofKernel.lean}
\item \nolinkurl{lean/Recombination/CarrierExactification/WitnessData.lean}
\item \nolinkurl{lean/Recombination/CarrierExactification/ProofGaps.lean}
\end{itemize}

The local target \texttt{Recombination.CarrierExactification} compiles, while the root
\nolinkurl{lean/Recombination.lean} remains intentionally unwired in the current worktree.
That layout makes the formal layer a local theorem module rather than a claim that the full
Lean root has already been reorganized around ticket31.

\subsection{Definitions and predicates}

\nolinkurl{Definitions.lean} provides the typed vocabulary behind the manuscript's main
distinctions. \texttt{DesignatedSurface} fixes the seed, observable, generator, section,
and protocol types. \texttt{CarrierCandidate} packages a state space together with
validity, seed realization, readout, and transport relations. The definitional layer then
introduces the predicates that parallel the paper's C1--C7 vocabulary:
\texttt{ReadoutSufficient}, \texttt{TransportClosed}, \texttt{TransportExact},
\texttt{ReadoutTransportCompatible}, \texttt{SectionHonest},
\texttt{ProtocolHonest}, \texttt{ExpansionHonest}, and \texttt{CrossSeedSupported}.

Reduction is represented by \texttt{Reduction}, \texttt{ReducibleTo}, and
\texttt{NonReducibleAmong}. The Lean layer also introduces explicit status types for scope,
law, compression, protocol indexing, quadratics, proto-events, and gate verdicts. This is
how the formal layer preserves the same distinctions as the main text: lawful witness
versus public shadow, scoped protocol-indexed evidence versus full exactification, and
non-faithful compression versus faithful compression.

\subsection{Criterion lemmas}

\nolinkurl{Statements.lean} and \nolinkurl{ProofKernel.lean} together provide the theorem
kernel referenced in Section~12. The kernel proves identity and composition for reductions,
reflexive and transitive reducibility, and preservation lemmas for validity, readout, and
transport under reduction. On top of that interface, it proves the main evidence-extraction
lemmas from both full exactification evidence and protocol-indexed evidence.

The theorem layer also encodes the gate's conservative implications. Full exactification
evidence implies scoped positivity, transport closure and exactness, section honesty,
protocol honesty, expansion honesty, cross-seed support, and non-reducibility. The theorem
\texttt{carrier\_exactified\_of\_evidence} packages this into the abstract exactification
criterion. The protocol-indexed layer proves that protocol-indexed evidence yields a scoped
positive carrier while still blocking any nonindexed promotion. Additional barrier lemmas
show that fully exactified conclusions do not follow when cross-seed, section, or protocol
barriers remain in force. Finally, \texttt{CompressionContext} supplies the abstract
faithfulness implication used to formalize the witness/public-shadow split at theorem level.

These lemmas formalize the abstract logic of the gate. They do not prove that the ticket31
artifact surface instantiates their premises. That empirical bridge remains external and is
made explicit again below.

\subsection{Witness-data bridge}

\nolinkurl{WitnessData.lean} provides the Lean-readable mirror of the final bounded verdict.
Its central structures are \texttt{CarrierLabel}, \texttt{ClaimPermissions}, and
\texttt{Ticket31WitnessData}. The instance \texttt{ticket31FinalWitnessData} mirrors the
paper's stable final rows through a short fact record:
\begin{itemize}[leftmargin=*]
\item final verdict \texttt{protocolIndexedCarrierOnly},
\item coarse public row \nolinkurl{intrinsic_naked_interface_coarse} with
\texttt{signatureOnly},
\item lawful witness
\nolinkurl{intrinsic_naked_interface_witness_rewrite_effective_operator} with
\texttt{lawful},
\item scoped survivor
\nolinkurl{intrinsic_naked_interface_witness_autonomous_protocol_lift} at
\texttt{flagshipOnly} scope with \texttt{protocolIndexedOnly} indexing, and
\item compression status \texttt{nonFaithful}.
\end{itemize}

The bridge also mirrors the scoped secondary rows and the permission boundary. Scoped
quadratic carrying and scoped proto-events are recorded as positive. The bridge keeps the
following broader permissions false:
\begin{itemize}[leftmargin=*]
\item public-surface law,
\item cross-seed carrier,
\item full-core carrier,
\item global quadratic carrying,
\item general event geometry, and
\item the excluded broader bridge row.
\end{itemize}
The bridge JSON \nolinkurl{results/derived/ticket31_lean_witness_bridge.json} records
exactly which facts were bridged into Lean and which remain external.

This is the key trust-boundary sentence for the appendix: the witness-data bridge is a
\emph{typed mirror} of the final vocabulary and permissions. It is not a proof that the
JSON artifacts, cargo-backed steps, or PICA-backed steps are correct.

\subsection{Proof gaps}

\nolinkurl{ProofGaps.lean} keeps the remaining external assumptions visible. In the current
snapshot there is no active \texttt{sorry} under the
\texttt{Recombination.CarrierExactification} namespace, but several explicit bridge gaps
remain. The most important for the manuscript are:
\begin{itemize}[leftmargin=*]
\item \texttt{GAP\_C1\_TO\_C7\_EMPIRICAL\_BRIDGE}
\item \texttt{GAP\_T19\_GATE\_SOUNDNESS\_FROM\_JSON}
\item \texttt{GAP\_C3\_NONREDUCIBILITY\_EMPIRICAL\_BRIDGE}
\item \texttt{GAP\_L03\_JSON\_EXTRACTION\_TRUST}
\item \texttt{GAP\_L03\_PICA\_RUNNER\_TRUST}
\end{itemize}

Two additional gap families matter for interpretation even though they are secondary to the
paper's main verdict: the semantic non-upgrade gap for protocol-indexed evidence, and the
missing higher-order bridge from the scoped quadratic/proto-event rows to a richer formal
interface. These gap names are useful because they make the trust boundary specific. The
external assumptions still include the correctness of the ticket31 JSON artifacts, the
correctness of the extraction scripts, the correctness of cargo/PICA and phase-runner
execution, the empirical C1--C7 and C3 audits, and the soundness of the final-gate bridge
from the artifact surface to the Lean witness-data mirror.

That boundary is not a defect in the appendix. It is the formal statement of what the paper
does and does not claim. Lean formalizes the vocabulary and abstract implication structure;
the witness-data layer mirrors the final bounded verdict; the empirical run remains outside
Lean and is supported instead by the canonical pipeline and asset index.
 \section{Extended Tables}
\label{app:extended-tables}

This appendix collects tables that support the main text but are too detailed to place in
the main flow. Each table adds detail beyond the corresponding main-text summary. None of
the tables changes the paper's claim boundary: the final verdict remains
\texttt{protocol\_indexed\_carrier\_only}, the lawful witness and public shadow remain
distinct, the nearby positives remain real scope limitations, and the scoped secondary
rows remain scoped.

Tables~\ref{tab:appendix-candidate-reduction-detail},
\ref{tab:appendix-seed-scope-contract-detail},
\ref{tab:appendix-t11m-detail}, \ref{tab:appendix-falsifier-detail}, and
\ref{tab:appendix-secondary-structure-detail} expand the candidate taxonomy, scope
contracts, strengthening lane, falsifier panels, and scoped secondary structures that are
summarized in the main text.

\begin{table}[p]
\centering
\scriptsize
\caption{Candidate census and reduction-map detail. This table expands the main taxonomy by
listing the concrete ticket31 candidate lanes and their reduction classification.}
\label{tab:appendix-candidate-reduction-detail}
\begin{tabularx}{\linewidth}{@{}>{\raggedright\arraybackslash}p{0.25\linewidth}>{\raggedright\arraybackslash}p{0.16\linewidth}>{\raggedright\arraybackslash}p{0.16\linewidth}>{\raggedright\arraybackslash}p{0.15\linewidth}>{\raggedright\arraybackslash}X@{}}
\toprule
Candidate id & Candidate class & Instantiation & Reduction classification & Appendix interpretation \\
\midrule
\nolinkurl{current_predictive_descriptor_baseline} & descriptor partition & instantiated & descriptor/process & baseline quotient layer only; no live carrier advance \\
\nolinkurl{branch_recombination_descriptor_baseline} & descriptor partition & instantiated & descriptor/process & branch/recombination baseline only \\
\nolinkurl{minimal_behavioral_carrier_baseline} & descriptor partition & instantiated & descriptor/process & \(\Lambda^\ast\)-style minimal baseline, not a promoted carrier \\
\nolinkurl{relative_cycle_process_execution_family} & process object family & instantiated & descriptor/process & runtime family retained as process evidence rather than carrier law \\
\nolinkurl{local_chart_intrinsic_sidebranch} & chart-only intrinsic & reference only & chart-only & local chart fit without global carrier status \\
\nolinkurl{intrinsic_naked_interface} & intrinsic carrier candidate & instantiated & live candidate & primary live-surface candidate leading to the later witness/public split \\
\nolinkurl{sectioned_package_interface} & chart plus section family & instantiated & unresolved & section-mediated lane; later remains chart-only rather than a landed carrier \\
\nolinkurl{bundle_section_fiber_interface} & section-fiber carrier candidate & instantiated & live candidate & secondary carrier-style lane, but more mediated than the intrinsic branch \\
\nolinkurl{protocol_indexed_interface} & protocol-indexed object & reference only & unresolved & early protocol-indexed lane fails to instantiate on the live surface; later scoped lift is a different object \\
\nolinkurl{candidate_extension_new_surface_family} & new-surface candidate & prospective & unresolved & prospective adjacent-surface lane only; later no exactification \\
\bottomrule
\end{tabularx}
\end{table}

\begin{table}[p]
\centering
\scriptsize
\caption{Seed-level scope-contract detail. This table expands the corpus table by showing
the T12.A/T12.B interpretation row by row.}
\label{tab:appendix-seed-scope-contract-detail}
\begin{tabularx}{\linewidth}{@{}>{\raggedright\arraybackslash}p{0.26\linewidth}>{\raggedright\arraybackslash}p{0.12\linewidth}>{\raggedright\arraybackslash}p{0.18\linewidth}>{\raggedright\arraybackslash}p{0.14\linewidth}>{\raggedright\arraybackslash}X@{}}
\toprule
Stable id & Role & T12.A reading & T12.B contract & Carrier interpretation \\
\midrule
\nolinkurl{t27.flagship.reference_positive} & flagship & anchor source row & anchor source row & ordinary flagship evidence \\
\nolinkurl{t27.flagship.completion_partner} & flagship & anchor source row & anchor source row & ordinary flagship evidence \\
\nolinkurl{t27.nearby.n3_d1_w1of3_2of3} & nearby positive & support-only alignment & real limit & blocks cross-seed/full-core promotion \\
\nolinkurl{t27.nearby.n4_d1_w1of3_2of3} & nearby positive & support-only alignment & real limit & blocks cross-seed/full-core promotion \\
\nolinkurl{t27.nearby.n5_d1_w1of2_1of2} & nearby positive & support-only alignment & real limit & blocks cross-seed/full-core promotion \\
\nolinkurl{t27.nearby.n5_d2_w1of3_2of3} & nearby positive & support-only alignment & real limit & blocks cross-seed/full-core promotion \\
\nolinkurl{t27.control.null_baseline} & control & control consistency row & null\_by\_design & not a transfer failure and not a positive carrier row \\
\nolinkurl{t27.control.memory_only} & control & control consistency row & null\_by\_design & not a transfer failure and not a positive carrier row \\
\nolinkurl{t27.control.artifact} & paired control & paired logical panel & null\_by\_design & remains separate from ordinary transfer counts \\
\nolinkurl{t27.control.boundary_collapse} & control & control consistency row & null\_by\_design & not a transfer failure and not a positive carrier row \\
\nolinkurl{t27.support.erasure_recovery_anchor} & support only & support-only row & positive\_transfer & positive only as support-only transfer, not exact readout or law transfer \\
\bottomrule
\end{tabularx}
\end{table}

\begin{table}[p]
\centering
\scriptsize
\caption{T11m strengthening-lane detail. This table expands the main strengthening summary
with the exact artifact row used at each step.}
\label{tab:appendix-t11m-detail}
\begin{tabularx}{\linewidth}{@{}>{\raggedright\arraybackslash}p{0.12\linewidth}>{\raggedright\arraybackslash}p{0.22\linewidth}>{\raggedright\arraybackslash}p{0.19\linewidth}>{\raggedright\arraybackslash}p{0.17\linewidth}>{\raggedright\arraybackslash}X@{}}
\toprule
Step & Artifact & Positive finding & Retained limit & Use in the paper \\
\midrule
T11m cutover & \nolinkurl{results/derived/ticket31_public_carrier_cutover_decision.json} & downstream basis set to \nolinkurl{witness_site_public_shadow} & no public-surface law upgrade & fixes the two-level operational basis \\
T11m-i & \nolinkurl{results/derived/ticket31_autonomous_protocol_surface_summary.json} & autonomous lifted object built; \texttt{indexed\_only}; 1776 states / 1776 readout entries / 6 transports on the lifted flagship patch & blocked stable ids and paired endpoints remain outside the lifted audited patch & establishes the scoped lifted surface \\
T11m-ii & \nolinkurl{results/derived/ticket31_lifted_loop_action_summary.json} & future-only separated pairs present & loop-action asymmetry remains \texttt{not\_found} & keeps future-only separation distinct from a stronger loop-action claim \\
T11m-iii & \nolinkurl{results/derived/ticket31_observable_enrichment_summary.json} & predictive witnesses salvaged; 2 families; 30 pair witnesses & control-side rows remain source-limited & adds flagship predictive-witness support without widening scope \\
Intersection check & \nolinkurl{results/derived/ticket31_predictive_witness_intersection.json} & all 30 observable predictive pairs lie inside the 78 future-separated base-pair population & loop-action predictive population remains empty & confirms that the enrichment lane is future-only, not loop-action based \\
T11m-iii-a & \nolinkurl{results/derived/ticket31_lifted_section_protocol_gate.json} & section verdict \texttt{section\_robust}; protocol verdict \texttt{protocol\_robust} on the two flagship stable ids & inherited control and paired-endpoint gaps remain explicit & closes honesty on the lifted flagship patch \\
T11m-iii-b & \nolinkurl{results/derived/ticket31_joint_attack_matrix.json} & predictive front closed; protocol front closed; joint status \texttt{both\_fronts\_closed} & scope still \texttt{flagship\_only}; loop-action still \texttt{not\_found} & justifies entering the intrinsic promotion gate \\
T11m-v & \nolinkurl{results/derived/ticket31_intrinsic_promotion_gate.json} & survivor recorded as \texttt{protocol\_indexed\_only} with one surviving candidate & no basis upgrade; recommended basis stays \nolinkurl{witness_site_public_shadow} & fixes the strongest scoped carrier row used by T17/T19 \\
\bottomrule
\end{tabularx}
\end{table}

\begin{table}[p]
\centering
\scriptsize
\caption{Falsifier-panel detail for T13--T16. This table expands the main stress and honesty
summary without reopening stronger claims.}
\label{tab:appendix-falsifier-detail}
\begin{tabularx}{\linewidth}{@{}>{\raggedright\arraybackslash}p{0.13\linewidth}>{\raggedright\arraybackslash}p{0.22\linewidth}>{\raggedright\arraybackslash}p{0.17\linewidth}>{\raggedright\arraybackslash}p{0.18\linewidth}>{\raggedright\arraybackslash}X@{}}
\toprule
Panel & Artifact & Branch decision & Preserved row & Blocking outcome \\
\midrule
T13 denominator stress & \nolinkurl{results/derived/ticket31_denominator_stress_summary.json} & fragile & lawful witness survives; falsifier verdict keeps the licensed law surface intact & coarse-public weakness, support-only overpromotion, and full-core overclaim remain blocked \\
T14 surface expansion & \nolinkurl{results/derived/ticket31_surface_expansion_summary.json} & fragile & lawful witness survives branch-consistent expansion checks & no adjacent-surface upgrade and no broader transfer claim \\
T15 section/protocol honesty & \nolinkurl{results/derived/ticket31_section_protocol_gate.json} & site-not-probe honesty fails & licensed basis and primary witness survive honesty in scoped form & global/full claim collapses; paired panel remains separate; route filtering unavailable \\
T16 adjacent-surface search & \nolinkurl{results/derived/ticket31_new_surface_exactification_summary.json} & no exactification & existing scoped basis remains the only positive lane & no new surface upgrades the verdict \\
\bottomrule
\end{tabularx}
\end{table}

\begin{table}[p]
\centering
\scriptsize
\caption{Scoped secondary-structure detail for T17 and T18. All positive rows remain tied to
the flagship-only, protocol-indexed surviving carrier.}
\label{tab:appendix-secondary-structure-detail}
\begin{tabularx}{\linewidth}{@{}>{\raggedright\arraybackslash}p{0.26\linewidth}>{\raggedright\arraybackslash}p{0.12\linewidth}>{\raggedright\arraybackslash}p{0.15\linewidth}>{\raggedright\arraybackslash}p{0.16\linewidth}>{\raggedright\arraybackslash}X@{}}
\toprule
Item & Status & Scope & Source artifact & Appendix interpretation \\
\midrule
Carried surface (\nolinkurl{intrinsic_naked_interface_witness_autonomous_protocol_lift}) & scoped carrier & flagship-only, protocol-indexed & \nolinkurl{results/derived/ticket31_intrinsic_promotion_gate.json} & the T17/T18 rows attach only to the survived scoped carrier \\
\nolinkurl{q_readout_l2_self} & carried & flagship-only, protocol-indexed & \nolinkurl{results/derived/ticket31_quadratic_carrying_summary.json} & claim-bearing quadratic family retained on the scoped carrier \\
\nolinkurl{q_readout_gap} & carried & flagship-only, protocol-indexed & \nolinkurl{results/derived/ticket31_quadratic_carrying_summary.json} & claim-bearing quadratic family retained on the scoped carrier \\
\nolinkurl{q_coarsening_consistency} & carried & flagship-only, protocol-indexed & \nolinkurl{results/derived/ticket31_quadratic_carrying_summary.json} & claim-bearing quadratic family retained on the scoped carrier \\
\nolinkurl{b_aligned_readout_inner_product} & comparison only & scoped comparison & \nolinkurl{results/derived/ticket31_quadratic_carrying_summary.json} & comparison row only; not promotable beyond the scoped carrier \\
\nolinkurl{flagship_quadratic_levelset_event} & present & flagship-only, protocol-indexed & \nolinkurl{results/derived/ticket31_proto_event_pilot_summary.json} & scoped proto-event/test family supported on ordinary flagship rows only \\
\nolinkurl{flagship_transition_stability_test} & present & flagship-only, protocol-indexed & \nolinkurl{results/derived/ticket31_proto_event_pilot_summary.json} & scoped transition/test family with protocol indexing required \\
\nolinkurl{flagship_protocol_indexed_readout_agreement_event} & absent & flagship-only, protocol-indexed & \nolinkurl{results/derived/ticket31_proto_event_pilot_summary.json} & absence of a third family blocks any broader event-geometry reading \\
\bottomrule
\end{tabularx}
\end{table}
 
\clearpage
\bibliographystyle{plainnat}
\bibliography{references}

\end{document}
